% =============================================================================
% BỘ CÔNG THƯƠNG - TRƯỜNG ĐẠI HỌC CÔNG NGHIỆP TP.HCM
% ĐỒ ÁN MÔN: CÔNG NGHỆ MỚI TRONG PHÁT TRIỂN ỨNG DỤNG CNTT
% ĐỀ TÀI: XÂY DỰNG ỨNG DỤNG NHẮN TIN ZALO CLONE
% =============================================================================
\documentclass[12pt,a4paper]{report}

% ---------------------- Ngôn ngữ & Font (LuaLaTeX) ------------
\usepackage{fontspec}
\setmainfont{DejaVu Serif}
\setsansfont{DejaVu Sans}
\setmonofont{DejaVu Sans Mono}
\usepackage{polyglossia}
\setmainlanguage{vietnamese}

% ---------------------- Định dạng trang ----------------------
\usepackage[top=2.5cm, bottom=2.5cm, left=3.5cm, right=2cm]{geometry}
\usepackage{setspace}
\onehalfspacing                % Dãn dòng 1.5

% ---------------------- Toán & Ký hiệu ----------------------
\usepackage{amsmath,amssymb}
\usepackage{textcomp}

% ---------------------- Đồ họa & Màu -------------------------
\usepackage{graphicx}
\usepackage[dvipsnames]{xcolor}
\usepackage{float}
\usepackage{subcaption}

% ---------------------- PlantUML (UML) -----------------------
% Biên dịch với: lualatex -shell-escape report.tex
% Cần tải plantuml.jar và set biến môi trường PLANTUML_JAR
\usepackage{plantuml}

% ---------------------- Liên kết -----------------------------
\usepackage[unicode, hidelinks]{hyperref}
\hypersetup{
    colorlinks=true,
    linkcolor=black,
    urlcolor=blue,
    citecolor=black
}

% ---------------------- Mã nguồn -----------------------------
\usepackage{listings}
\definecolor{codegreen}{rgb}{0,0.6,0}
\definecolor{codegray}{rgb}{0.5,0.5,0.5}
\definecolor{codepurple}{rgb}{0.58,0,0.82}
\definecolor{backcolour}{rgb}{0.95,0.95,0.92}

\lstdefinestyle{mystyle}{
    backgroundcolor=\color{backcolour},
    commentstyle=\color{codegreen},
    keywordstyle=\color{magenta},
    numberstyle=\tiny\color{codegray},
    stringstyle=\color{codepurple},
    basicstyle=\ttfamily\footnotesize,
    breakatwhitespace=false,
    breaklines=true,
    captionpos=b,
    keepspaces=true,
    numbers=left,
    numbersep=5pt,
    showspaces=false,
    showstringspaces=false,
    showtabs=false,
    tabsize=2,
    frame=single
}
\lstset{style=mystyle}

% ---------------------- Bảng biểu ----------------------------
\usepackage{booktabs}
\usepackage{longtable}
\usepackage{tabularx}
\usepackage{xltabular}
\usepackage{array}

% ---------------------- Header/Footer ------------------------
\usepackage{fancyhdr}
\pagestyle{fancy}
\fancyhf{}
\fancyhead[L]{}
\fancyhead[R]{\thepage}
\fancyfoot[C]{}

% ---------------------- Tùy chỉnh chapter --------------------
\usepackage{titlesec}
\titleformat{\chapter}[display]
  {\centering\normalfont\large\bfseries}
  {\MakeUppercase{\chaptertitlename}\ \thechapter}{0pt}{\normalfont\large}
\titlespacing*{\chapter}{0pt}{-10pt}{10pt}

% ---------------------- Tiêu đề section ----------------------
\titleformat{\section}{\normalfont\normalsize\bfseries}{\thesection}{1em}{}
\titleformat{\subsection}{\normalfont\normalsize\bfseries}{\thesubsection}{1em}{}
\titleformat{\subsubsection}{\normalfont\normalsize\bfseries}{\thesubsubsection}{1em}{}

% ---------------------- Custom commands ----------------------
\newcommand{\nhom}{16}
\newcommand{\hMot}{Thái Viết Phát}
\newcommand{\mssvMot}{22001265}
\newcommand{\hHai}{Nguyễn Ngọc Tài}
\newcommand{\mssvHai}{21000881}
\newcommand{\hBa}{Nguyễn Phi Thiên}
\newcommand{\mssvBa}{20000735}
\newcommand{\hBon}{Nguyễn Huy Hoàng}
\newcommand{\mssvBon}{20100521}
\newcommand{\tenDeTai}{Xây dựng ứng dụng nhắn tin Zalo Clone}

\newcommand{\blankpage}{\newpage\null\thispagestyle{empty}\newpage}

% =============================================================================
\title{\tenDeTai}
\author{}
\date{}

% =============================================================================
\begin{document}

% ===================== BÌA ===================================================
\begin{titlepage}
\begin{center}

% Logo & Header
\begin{figure}[t]
    \centering
    \includegraphics[width=0.9\textwidth]{logo.png}
\end{figure}

\textbf{BỘ CÔNG THƯƠNG} \\
\textbf{TRƯỜNG ĐẠI HỌC CÔNG NGHIỆP TP.HCM} \\
\textbf{---o0o---}

\vspace{2cm}

\textbf{\large ĐỒ ÁN MÔN HỌC} \\
\textbf{CÔNG NGHỆ MỚI TRONG PHÁT TRIỂN ỨNG DỤNG CNTT}

\vspace{2cm}

\textbf{\Large \tenDeTai}

\vspace{2cm}

\textbf{Nhóm \nhom}

\vspace{1cm}

\begin{tabular}{lcl}
    \hMot & --- & \mssvMot \\
    \hHai & --- & \mssvHai \\
    \hBa  & --- & \mssvBa \\
    \hBon & --- & \mssvBon \\
\end{tabular}

\vspace{2cm}

\textbf{TP. HỒ CHÍ MINH, 2026}

\end{center}
\end{titlepage}

% ===================== MỤC LỤC ===============================================
\tableofcontents
\thispagestyle{empty}

% ===================== DANH MỤC HÌNH =========================================
\newpage
\listoffigures
\addcontentsline{toc}{chapter}{DANH MỤC CÁC HÌNH VẼ}

% ===================== DANH MỤC BẢNG =========================================
\newpage
\listoftables
\addcontentsline{toc}{chapter}{DANH MỤC CÁC BẢNG BIỂU}

\newpage
\setcounter{page}{1}

% #############################################################################
% CHƯƠNG 1: GIỚI THIỆU
% #############################################################################
\chapter{GIỚI THIỆU}

\section{Tổng quan}

Trong thời đại số hóa hiện nay, các ứng dụng nhắn tin đã trở thành một phần không thể thiếu trong cuộc sống hàng ngày của hàng tỷ người dùng trên toàn thế giới. Từ việc kết nối bạn bè, gia đình đến trao đổi công việc, các ứng dụng như Zalo, Messenger, WhatsApp đã thay đổi căn bản cách con người giao tiếp.

Zalo là một trong những ứng dụng nhắn tin phổ biến nhất tại Việt Nam với hơn 70 triệu người dùng, cung cấp đa dạng các tính năng từ nhắn tin văn bản, gọi điện thoại và video, chia sẻ file, đến các tính năng mạng xã hội như feed và AI chat.

Đồ án này nhằm mục đích xây dựng một ứng dụng nhắn tin tương tự Zalo (Zalo Clone) với đầy đủ các tính năng cốt lõi, sử dụng công nghệ hiện đại: Spring Boot cho backend, React cho frontend, WebSocket cho giao tiếp thời gian thực, và triển khai trên nền tảng đám mây.

\section{Mục tiêu đề tài}

Đề tài được thực hiện với các mục tiêu chính sau:

\begin{itemize}
    \item \textbf{Xây dựng hệ thống backend} hoàn chỉnh với kiến trúc RESTful API, quản lý người dùng, xác thực JWT, và giao tiếp WebSocket thời gian thực.
    \item \textbf{Xây dựng giao diện frontend} hiện đại với React 19, hỗ trợ real-time messaging, quản lý trạng thái tập trung với Zustand.
    \item \textbf{Hệ thống nhắn tin 1-1 và nhóm} với hỗ trợ đa phương tiện (hình ảnh, video, âm thanh, file).
    \item \textbf{Tính năng xã hội}: kết bạn, chặn người dùng, báo cáo vi phạm.
    \item \textbf{Tích hợp AI Chatbot} sử dụng Groq API (LLaMA 3.1).
    \item \textbf{Triển khai lên nền tảng đám mây}: Railway (backend), Vercel (frontend), Cloudflare R2 (lưu trữ file).
\end{itemize}

\section{Phạm vi đề tài}

Đồ án tập trung vào phát triển một ứng dụng nhắn tin web với các phạm vi sau:

\begin{itemize}
    \item \textbf{Nền tảng}: Web application (responsive), hỗ trợ trình duyệt hiện đại.
    \item \textbf{Backend}: Java 21, Spring Boot 3.4.2, MySQL, Redis, WebSocket.
    \item \textbf{Frontend}: React 19, Vite, Tailwind CSS, Zustand.
    \item \textbf{Triển khai}: Railway (backend), Vercel (frontend), Cloudflare R2 (file storage).
    \item \textbf{Không bao gồm}: Ứng dụng mobile native, gọi điện thoại/video real-time (Call Session đã được ghi nhận log), end-to-end encryption.
\end{itemize}

\section{Mô tả yêu cầu chức năng}

\subsection{Phân loại người dùng}

Hệ thống có ba nhóm tác nhân chính:

\begin{itemize}
    \item \textbf{Khách (Guest)}: Người dùng chưa đăng nhập. Có thể đăng ký tài khoản mới.
    \item \textbf{Người dùng (User)}: Đã đăng nhập. Có thể nhắn tin, kết bạn, tạo nhóm, sử dụng AI chat, đăng bài.
    \item \textbf{Quản trị viên (Admin)}: Quản lý người dùng, nhóm, xử lý báo cáo, xem thống kê.
\end{itemize}

\subsection{Danh sách chức năng chính}

\begin{enumerate}
    \item \textbf{Xác thực và bảo mật}: đăng ký, đăng nhập, đăng xuất, refresh token, quên mật khẩu, xác thực email.
    \item \textbf{Quản lý hồ sơ}: xem và cập nhật thông tin cá nhân, avatar, tìm kiếm người dùng.
    \item \textbf{Kết bạn}: gửi/chấp nhận/từ chối lời mời kết bạn, danh sách bạn bè.
    \item \textbf{Nhắn tin 1-1}: gửi tin nhắn văn bản, hình ảnh, video, file, âm thanh; đánh dấu đã xem; thu hồi và xóa tin nhắn.
    \item \textbf{Nhắn tin nhóm}: tạo nhóm, thêm/xóa thành viên, phân quyền quản trị nhóm.
    \item \textbf{Reaction}: thả cảm xúc (emoji) cho tin nhắn 1-1 và tin nhắn nhóm.
    \item \textbf{AI Chatbot}: trò chuyện với AI (Groq/LLaMA 3.1) trong giao diện chat riêng.
    \item \textbf{Chặn và báo cáo}: chặn người dùng, báo cáo vi phạm, admin xử lý báo cáo.
    \item \textbf{Call Session}: ghi nhận lịch sử cuộc gọi thoại/video.
    \item \textbf{Quản trị}: thống kê dashboard, quản lý người dùng/nhóm, audit log.
\end{enumerate}

% #############################################################################
% CHƯƠNG 2: CƠ SỞ LÝ THUYẾT
% #############################################################################
\chapter{CƠ SỞ LÝ THUYẾT}

\section{Tổng quan công nghệ}

Hệ thống được xây dựng dựa trên mô hình kiến trúc client-server với các công nghệ hiện đại:

\begin{itemize}
    \item \textbf{Backend}: Spring Boot 3.4.2 (Java 21) --- framework phát triển ứng dụng Java mạnh mẽ, cung cấp dependency injection, bảo mật tích hợp, và hỗ trợ WebSocket native.
    \item \textbf{Frontend}: React 19 kết hợp Vite --- thư viện xây dựng giao diện người dùng phổ biến, với Vite là công cụ build nhanh dựa trên ESBuild.
    \item \textbf{Cơ sở dữ liệu}: MySQL --- hệ quản trị cơ sở dữ liệu quan hệ phổ biến, quản lý schema với Flyway migrations.
    \item \textbf{Real-time}: STOMP over WebSocket với SockJS fallback --- giao thức nhắn tin publish/subscribe cho ứng dụng thời gian thực.
    \item \textbf{Container hóa \& Triển khai}: Docker, Railway (backend), Vercel (frontend), Cloudflare R2 (file storage).
\end{itemize}

\section{Spring Boot}

Spring Boot là một framework phát triển ứng dụng Java dựa trên nền tảng Spring. Nó cung cấp cấu hình tự động (auto-configuration), máy chủ embedded (Tomcat, Jetty), và hệ sinh thái phong phú các starter cho phép phát triển nhanh các ứng dụng microservice. Trong đồ án này, Spring Boot được sử dụng với:

\begin{itemize}
    \item \textbf{Spring Data JPA}: Tầng truy cập dữ liệu, ánh xạ object-relational (ORM).
    \item \textbf{Spring Security + JWT}: Xác thực và phân quyền với JSON Web Token.
    \item \textbf{Spring WebSocket}: Giao tiếp thời gian thực với STOMP protocol.
    \item \textbf{Spring Validation}: Kiểm tra dữ liệu đầu vào với Jakarta Validation.
    \item \textbf{Flyway}: Quản lý phiên bản cơ sở dữ liệu.
\end{itemize}

\section{React và Hệ sinh thái}

React là thư viện JavaScript mã nguồn mở do Facebook phát triển, dùng để xây dựng giao diện người dùng dựa trên kiến trúc component. Đồ án sử dụng React 19 với:

\begin{itemize}
    \item \textbf{Vite}: Công cụ build thế hệ mới, nhanh hơn Webpack nhờ ESBuild và HMR (Hot Module Replacement) hiệu quả.
    \item \textbf{Zustand 5}: Thư viện quản lý state đơn giản, nhẹ, dựa trên hooks, thay thế Redux cho các ứng dụng vừa và nhỏ.
    \item \textbf{Tailwind CSS 3}: Framework CSS utility-first, xây dựng giao diện nhanh với các lớp CSS có sẵn.
    \item \textbf{Axios}: HTTP client cho trình duyệt, hỗ trợ interceptor để tự động đính kèm JWT và refresh token.
    \item \textbf{@stomp/stompjs + sockjs-client}: STOMP client cho WebSocket với fallback khi WebSocket không khả dụng.
\end{itemize}

\section{WebSocket và STOMP}

\textbf{WebSocket} là giao thức truyền thông hai chiều qua một kết nối TCP duy nhất, khác với HTTP truyền thống là request-response. WebSocket cho phép server chủ động gửi dữ liệu đến client mà không cần client yêu cầu --- rất phù hợp cho ứng dụng chat.

\textbf{STOMP} (Streaming Text Oriented Messaging Protocol) là giao thức nhắn tin đơn giản, hoạt động trên WebSocket, cung cấp cơ chế publish/subscribe. Client có thể subscribe vào một "destination" để nhận message, và gửi message đến destination khác.

\begin{figure}[H]
    \centering
    \begin{plantuml}
        @startuml
        !theme plain
        actor "Client A" as A
        actor "Client B" as B
        participant "STOMP Broker" as Broker
        database "Backend" as BE

        A -> Broker: SUBSCRIBE /topic/chat/123
        B -> Broker: SUBSCRIBE /topic/chat/123
        B -> BE: SEND /app/chat/123/sendMessage
        BE -> Broker: MESSAGE /topic/chat/123
        Broker -> A: MESSAGE (new message)
        Broker -> B: MESSAGE (new message)
        @enduml
    \end{plantuml}
    \caption{Mô hình giao tiếp WebSocket/STOMP}
\end{figure}

\section{JWT (JSON Web Token)}

JWT là chuẩn mở (RFC 7519) định nghĩa cách truyền thông tin an toàn giữa các bên dưới dạng JSON object. Token gồm ba phần: header, payload, signature. Trong đồ án:

\begin{itemize}
    \item \textbf{Access Token}: Thời hạn 24 giờ, chứa email, userId, role, tokenVersion. Được gửi trong header Authorization cho mọi request.
    \item \textbf{Refresh Token}: Thời hạn 7 ngày, dùng để lấy access token mới khi hết hạn.
    \item \textbf{Token Version}: Cơ chế invalidate toàn bộ session cũ khi đăng nhập từ thiết bị mới, tăng tokenVersion trong DB và so sánh với version trong JWT.
\end{itemize}

\section{Cloud Computing --- Railway và Cloudflare R2}

\textbf{Railway} là nền tảng cloud PaaS (Platform as a Service) hỗ trợ triển khai ứng dụng với nhiều ngôn ngữ. Dự án sử dụng Railway để:

\begin{itemize}
    \item Triển khai backend Spring Boot dưới dạng Docker container.
    \item Cung cấp MySQL addon (cùng internal network, latency < 1ms).
    \item Tự động health check và restart khi gặp lỗi.
\end{itemize}

\textbf{Cloudflare R2} là dịch vụ lưu trữ object tương thích S3, cung cấp 10GB miễn phí và không tính phí egress. Dự án sử dụng R2 để lưu trữ file tải lên (hình ảnh, video, file đính kèm) thông qua AWS SDK v2.

\section{Groq AI --- LLaMA 3.1}

Groq là nền tảng inference AI với tốc độ cao, sử dụng LPU (Language Processing Unit) riêng. Dự án tích hợp Groq API với model LLaMA 3.1 8B, tương thích OpenAI API format, thông qua Spring AI. Chatbot AI có thể trả lời câu hỏi, hỗ trợ công việc trong giao diện chat riêng biệt.

% #############################################################################
% CHƯƠNG 3: PHÂN TÍCH VÀ THIẾT KẾ
% #############################################################################
\chapter{PHÂN TÍCH VÀ THIẾT KẾ}

\section{Phân tích yêu cầu bằng UML}

\subsection{Use case tổng quát}

\begin{figure}[H]
    \centering
    \begin{plantuml}
        @startuml
        !theme plain
        skinparam usecaseFontSize 12
        skinparam actorFontSize 12
        skinparam defaultFontSize 12

        actor "Khách\n(Guest)" as G
        actor "Người dùng\n(User)" as U
        actor "Quản trị viên\n(Admin)" as A

        (Đăng ký) as UC_REG
        (Đăng nhập) as UC_LOGIN
        (Quản lý hồ sơ) as UC_PROFILE
        (Nhắn tin 1-1) as UC_CHAT
        (Nhắn tin nhóm) as UC_GROUP
        (Kết bạn) as UC_FRIEND
        (Thả cảm xúc) as UC_REACT
        (AI Chat) as UC_AI
        (Ghi log cuộc gọi) as UC_CALL
        (Chặn người dùng) as UC_BLOCK
        (Báo cáo vi phạm) as UC_REPORT
        (Xử lý báo cáo) as UC_RESOLVE
        (Quản lý người dùng) as UC_ADMIN_USER
        (Xem thống kê) as UC_STATS

        G --> UC_REG
        G --> UC_LOGIN
        U --> UC_PROFILE
        U --> UC_CHAT
        U --> UC_GROUP
        U --> UC_FRIEND
        U --> UC_REACT
        U --> UC_AI
        U --> UC_CALL
        U --> UC_BLOCK
        U --> UC_REPORT
        A --> UC_BLOCK
        A --> UC_REPORT
        A --> UC_RESOLVE
        A --> UC_ADMIN_USER
        A --> UC_STATS
        @enduml
    \end{plantuml}
    \caption{Use case tổng quát của hệ thống}
\end{figure}

\subsection{Danh sách tác nhân và mô tả}

\begin{table}[H]
    \centering
    \caption{Danh sách tác nhân và mô tả}
    \begin{tabularx}{\textwidth}{|c|X|}
        \hline
        \textbf{Tác nhân} & \textbf{Mô tả} \\
        \hline
        Khách (Guest) & Người dùng chưa đăng nhập, có thể đăng ký tài khoản mới \\
        \hline
        Người dùng (User) & Thành viên đã đăng nhập, sử dụng toàn bộ tính năng chat và xã hội \\
        \hline
        Quản trị viên (Admin) & Người dùng có quyền quản trị, quản lý hệ thống và xử lý vi phạm \\
        \hline
    \end{tabularx}
\end{table}

\subsection{Danh sách các tình huống sử dụng}

\begin{longtable}{|c|p{11cm}|}
    \hline
    \textbf{ID} & \textbf{Tên Use case} \\
    \hline
    UC01 & Đăng ký tài khoản \\
    \hline
    UC02 & Đăng nhập \\
    \hline
    UC03 & Cập nhật hồ sơ cá nhân \\
    \hline
    UC04 & Tìm kiếm người dùng \\
    \hline
    UC05 & Gửi lời mời kết bạn \\
    \hline
    UC06 & Chấp nhận lời mời kết bạn \\
    \hline
    UC07 & Nhắn tin văn bản 1-1 \\
    \hline
    UC08 & Gửi file đa phương tiện \\
    \hline
    UC09 & Tạo nhóm chat \\
    \hline
    UC10 & Thêm thành viên vào nhóm \\
    \hline
    UC11 & Thả cảm xúc cho tin nhắn \\
    \hline
    UC12 & Trò chuyện với AI \\
    \hline
    UC13 & Chặn người dùng \\
    \hline
    UC14 & Báo cáo người dùng \\
    \hline
    UC15 & Xử lý báo cáo (Admin) \\
    \hline
    UC16 & Xem thống kê hệ thống (Admin) \\
    \hline
    UC17 & Ghi log cuộc gọi \\
    \hline
    \caption{Danh sách các tình huống sử dụng trong hệ thống}
\end{longtable}

\subsection{Đặc tả use case}

\subsubsection{UC01 --- Đăng ký tài khoản}

\renewcommand{\arraystretch}{1.4}
\begin{xltabular}{\textwidth}{|X|X|}
\hline
\multicolumn{2}{|p{0.97\textwidth}|}{\textbf{Tên usecase:} Đăng ký.} \\
\hline
\multicolumn{2}{|p{0.97\textwidth}|}{\textbf{Actor:} Người dùng.} \\
\hline
\multicolumn{2}{|p{0.97\textwidth}|}{\textbf{Mô tả:} Cho phép người dùng đăng ký tài khoản.} \\
\hline
\multicolumn{2}{|p{0.97\textwidth}|}{\textbf{Tiền điều kiện:} Không có.} \\
\hline
\multicolumn{2}{|p{0.97\textwidth}|}{\textbf{Hậu điều kiện:} Đăng ký thành công.} \\
\hline
\multicolumn{2}{|c|}{\textbf{Luồng xử lý chính (Basic flow)}} \\
\hline
\multicolumn{1}{|c|}{\textbf{Actor}} & \multicolumn{1}{c|}{\textbf{System}} \\
\hline
1. Tại giao diện đăng nhập, chọn ``Đăng ký''. & \\
\hline
& 2. Hiển thị giao diện đăng ký yêu cầu nhập email, mật khẩu, xác nhận mật khẩu. \\
\hline
3. Nhập email, mật khẩu, xác nhận mật khẩu. & \\
\hline
4. Nhấn ``Đăng ký''. & \\
\hline
& 5. Kiểm tra email. \\
\hline
& 6. Chuyển sang giao diện nhập mã OTP. \\
\hline
7. Nhập mã OTP. & \\
\hline
8. Nhấn ``Xác nhận''. & \\
\hline
& 9. Kiểm tra mã OTP. \\
\hline
& 10. Chuyển sang giao diện nhập tên và mật khẩu. \\
\hline
11. Nhập thông tin tên và mật khẩu. & \\
\hline
12. Nhấn ``Hoàn tất đăng ký''. & \\
\hline
& 13. Lưu thông tin user vào database và thông báo đăng ký thành công. \\
\hline
\multicolumn{2}{|p{0.97\textwidth}|}{\textbf{Luồng sự kiện thay thế (Alternative Flows)}} \\
\hline
\multicolumn{2}{|p{0.97\textwidth}|}{5.1 Hiển thị thông báo email đã được đăng ký, yêu cầu nhập lại.} \\
\hline
\multicolumn{2}{|p{0.97\textwidth}|}{5.2 Quay về bước 2.} \\
\hline
\multicolumn{2}{|p{0.97\textwidth}|}{9.1 Hiển thị thông báo lỗi xác thực OTP thất bại.} \\
\hline
\multicolumn{2}{|p{0.97\textwidth}|}{9.2 Quay về bước 5.} \\
\hline
\caption{Mô tả chức năng UC01}
\end{xltabular}
\renewcommand{\arraystretch}{1}

\begin{figure}[H]
    \centering
    \begin{plantuml}
        @startuml
        !theme plain
        title Đăng ký
        |Actor|
        start
        :1. Tại giao diện đăng nhập, chọn "Đăng ký";
        |System|
        :2. Hiển thị giao diện đăng ký\nyêu cầu nhập email, mật khẩu, xác nhận mật khẩu;
        |Actor|
        :3. Nhập email, mật khẩu, xác nhận mật khẩu;
        :4. Nhấn "Đăng ký";
        |System|
        if (5. Kiểm tra email) then (Đã tồn tại)
          #pink:5.1 Hiển thị thông báo email đã\nđược đăng ký, yêu cầu nhập lại;
          :5.2 Quay về bước 2;
          detach
        else (Chưa tồn tại)
          :6. Chuyển sang giao diện nhập mã OTP;
        endif
        |Actor|
        :7. Nhập mã OTP;
        :8. Nhấn "Xác nhận";
        |System|
        if (9. Kiểm tra mã OTP) then (Không hợp lệ)
          #pink:9.1 Hiển thị thông báo lỗi\nxác thực OTP thất bại;
          :9.2 Quay về bước 5;
          detach
        else (Hợp lệ)
          :10. Chuyển sang giao diện nhập tên và mật khẩu;
        endif
        |Actor|
        :11. Nhập thông tin tên và mật khẩu;
        :12. Nhấn "Hoàn tất đăng ký";
        |System|
        :13. Lưu thông tin user vào database\nvà thông báo đăng ký thành công;
        stop
        @enduml
    \end{plantuml}
    \caption{Activity diagram --- Đăng ký tài khoản}
\end{figure}

\begin{figure}[H]
    \centering
    \begin{plantuml}
        @startuml
        !theme plain
        actor "User" as U
        participant "GD_DangNhap\n(Đăng nhập)" as GD1
        participant "GD_DangKy\n(Đăng ký)" as GD2
        participant "GD_NhapOTP\n(Nhập OTP)" as GD3
        participant "GD_NhapThongTin\n(Nhập thông tin)" as GD4
        participant "Ctrl_DangKy\n(Controller)" as CTRL
        database "Database" as DB

        U -> GD1: 1. Tại giao diện đăng nhập, chọn "Đăng ký"
        GD1 -> GD2: 2. Hiển thị giao diện đăng ký\nyêu cầu nhập email, mật khẩu, xác nhận mật khẩu
        activate GD2
        GD2 --> U: 2.1 return

        U -> GD2: 3. Nhập email, mật khẩu, xác nhận mật khẩu
        U -> GD2: 4. Nhấn "Đăng ký"
        GD2 -> CTRL: 5. Kiểm tra email
        activate CTRL
        CTRL -> DB: existsByEmail(email): boolean
        activate DB
        DB --> CTRL: 5.1 return
        deactivate DB

        alt [Email đã tồn tại]
          CTRL --> GD2: 5.1.1 Thông báo email đã được đăng ký, yêu cầu nhập lại
          GD2 --> U: 5.2 Quay về bước 2
        else [Email chưa tồn tại]
          CTRL --> GD2: 6. Chuyển sang giao diện nhập mã OTP
          deactivate CTRL
          GD2 -> GD3: 6. Hiển thị GD_NhapOTP
          activate GD3
          GD3 --> U: return

          U -> GD3: 7. Nhập mã OTP
          U -> GD3: 8. Nhấn "Xác nhận"
          GD3 -> CTRL: 9. Kiểm tra mã OTP
          activate CTRL

          alt [OTP không hợp lệ]
            CTRL --> GD3: 9.1 Hiển thị thông báo lỗi xác thực OTP thất bại
            GD3 --> U: 9.2 Quay về bước 5
          else [OTP hợp lệ]
            CTRL --> GD3: 10. Chuyển sang giao diện nhập tên và mật khẩu
            deactivate CTRL
            GD3 -> GD4: 10. Hiển thị GD_NhapThongTin
            activate GD4
            GD4 --> U: return

            U -> GD4: 11. Nhập thông tin tên và mật khẩu
            U -> GD4: 12. Nhấn "Hoàn tất đăng ký"
            GD4 -> CTRL: 13. Lưu thông tin user
            activate CTRL
            CTRL -> DB: createUser(user): User
            activate DB
            DB --> CTRL: 13.1.1 return
            deactivate DB
            CTRL --> GD4: 13.2 return
            deactivate CTRL
            GD4 --> U: 13.3 Thông báo đăng ký thành công
          end
        end
        @enduml
    \end{plantuml}
    \caption{Sequence diagram --- Đăng ký tài khoản}
\end{figure}

\subsubsection{UC02 --- Đăng nhập}

\renewcommand{\arraystretch}{1.4}
\begin{xltabular}{\textwidth}{|X|X|}
\hline
\multicolumn{2}{|p{0.97\textwidth}|}{\textbf{Tên usecase:} Đăng nhập.} \\
\hline
\multicolumn{2}{|p{0.97\textwidth}|}{\textbf{Actor:} Người dùng.} \\
\hline
\multicolumn{2}{|p{0.97\textwidth}|}{\textbf{Mô tả:} Cho phép người dùng đăng nhập vào tài khoản của mình.} \\
\hline
\multicolumn{2}{|p{0.97\textwidth}|}{\textbf{Tiền điều kiện:} Người dùng đã có tài khoản.} \\
\hline
\multicolumn{2}{|p{0.97\textwidth}|}{\textbf{Hậu điều kiện:} Đăng nhập thành công.} \\
\hline
\multicolumn{2}{|c|}{\textbf{Luồng xử lý chính (Basic flow)}} \\
\hline
\multicolumn{1}{|c|}{\textbf{Actor}} & \multicolumn{1}{c|}{\textbf{System}} \\
\hline
1. Tại giao diện đăng nhập, nhập email và mật khẩu. & \\
\hline
2. Nhấn ``Đăng nhập''. & \\
\hline
& 3. Kiểm tra email tồn tại. \\
\hline
& 4. Kiểm tra mật khẩu. \\
\hline
& 5. Hiển thị thông báo đăng nhập thành công và lưu thông tin đăng nhập vào database. \\
\hline
\multicolumn{2}{|p{0.97\textwidth}|}{\textbf{Luồng sự kiện thay thế (Alternative Flows)}} \\
\hline
\multicolumn{2}{|p{0.97\textwidth}|}{3.1 Hiển thị thông báo email chưa được đăng kí.} \\
\hline
\multicolumn{2}{|p{0.97\textwidth}|}{4.1 Hiển thị thông báo mật khẩu không đúng.} \\
\hline
\multicolumn{2}{|p{0.97\textwidth}|}{4.2 Quay về bước 1.} \\
\hline
\caption{Mô tả chức năng UC02}
\end{xltabular}
\renewcommand{\arraystretch}{1}

\begin{figure}[H]
    \centering
    \begin{plantuml}
        @startuml
        !theme plain
        title Đăng nhập
        |Người dùng|
        start
        :1. Tại giao diện đăng nhập,\nnhập email và mật khẩu;
        :2. Nhấn "Đăng nhập";
        |System|
        if (3. Kiểm tra email tồn tại) then (Đã tồn tại)
          if (4. Kiểm tra mật khẩu) then (Hợp lệ)
            :5. Hiển thị thông báo đăng nhập thành công\nvà lưu thông tin đăng nhập vào database;
          else (Không hợp lệ)
            #pink:4.1 Hiển thị thông báo\nmật khẩu không đúng;
            :4.2 Quay về bước 1;
            detach
          endif
        else (Chưa tồn tại)
          #pink:3.1 Hiển thị thông báo\nemail chưa được đăng kí;
          detach
        endif
        stop
        @enduml
    \end{plantuml}
    \caption{Activity diagram --- Đăng nhập}
\end{figure}

\begin{figure}[H]
    \centering
    \begin{plantuml}
        @startuml
        !theme plain
        actor "Người dùng" as U
        participant "GD_DangNhap\n(Giao diện đăng nhập)" as GD
        participant "Ctrl_Login\n(Controller)" as CTRL
        participant "User\n(Entity/Model)" as UM
        database "LoginLog\n(Database)" as LL

        U -> GD: 1. Tại giao diện đăng nhập, nhập email và mật khẩu
        U -> GD: 2. Nhấn "Đăng nhập"
        GD -> CTRL: 2.1 Gửi yêu cầu đăng nhập (email, mật khẩu)
        activate CTRL

        CTRL -> UM: 3. existsByEmail(email): boolean
        activate UM
        UM --> CTRL: 3.1 return
        deactivate UM

        alt [Email chưa được đăng ký]
          CTRL --> GD: 3.1 Hiển thị thông báo email chưa được đăng kí
          GD --> U: Thông báo lỗi
        else [Email đã đăng ký]
          CTRL -> UM: 4. isPasswordValid(email, password): boolean
          activate UM
          UM --> CTRL: 4.1 return
          deactivate UM

          alt [Mật khẩu không đúng]
            CTRL --> GD: 4.1 Hiển thị thông báo mật khẩu không đúng
            GD --> U: 4.2 Quay về bước 1
          else [Mật khẩu đúng]
            CTRL -> LL: 5. save(userId, loginTime)
            activate LL
            LL --> CTRL: 5.1.1 return
            deactivate LL
            CTRL --> GD: 5.2 Hiển thị thông báo đăng nhập thành công
            deactivate CTRL
            GD --> U: 5.3 Kết quả đăng nhập thành công
          end
        end
        @enduml
    \end{plantuml}
    \caption{Sequence diagram --- Đăng nhập}
\end{figure}

\subsubsection{UC03 --- Cập nhật hồ sơ cá nhân}

\renewcommand{\arraystretch}{1.4}
\begin{xltabular}{\textwidth}{|X|X|}
\hline
\multicolumn{2}{|p{0.97\textwidth}|}{\textbf{Tên usecase:} Cập nhật hồ sơ cá nhân.} \\
\hline
\multicolumn{2}{|p{0.97\textwidth}|}{\textbf{Actor:} Người dùng.} \\
\hline
\multicolumn{2}{|p{0.97\textwidth}|}{\textbf{Mô tả:} Cho phép người dùng chỉnh sửa thông tin cá nhân như tên, ảnh đại diện, và các thông tin cơ bản khác.} \\
\hline
\multicolumn{2}{|p{0.97\textwidth}|}{\textbf{Tiền điều kiện:}} \\
\hline
\multicolumn{2}{|p{0.97\textwidth}|}{\quad $+$ Người dùng đã có tài khoản hợp lệ trên hệ thống.} \\
\hline
\multicolumn{2}{|p{0.97\textwidth}|}{\quad $+$ Người dùng đã đăng nhập vào hệ thống.} \\
\hline
\multicolumn{2}{|p{0.97\textwidth}|}{\quad $+$ Có quyền truy cập vào trang hồ sơ cá nhân.} \\
\hline
\multicolumn{2}{|p{0.97\textwidth}|}{\textbf{Hậu điều kiện:}} \\
\hline
\multicolumn{2}{|p{0.97\textwidth}|}{\quad $+$ \textbf{Thành công:} Thông tin cá nhân được cập nhật và lưu vào hệ thống.} \\
\hline
\multicolumn{2}{|p{0.97\textwidth}|}{\quad $+$ \textbf{Thất bại:} Dữ liệu không thay đổi, hiển thị thông báo lỗi phù hợp.} \\
\hline
\multicolumn{2}{|c|}{\textbf{Luồng sự kiện chính (Basic flows):}} \\
\hline
\multicolumn{1}{|c|}{\textbf{Actor}} & \multicolumn{1}{c|}{\textbf{System}} \\
\hline
1. Truy cập vào trang ``Hồ sơ cá nhân''. & \\
\hline
& 2. Hiển thị thông tin hiện tại (tên, ảnh, email,\ldots) \\
\hline
3. Chọn ``Chỉnh sửa sơ''. & \\
\hline
4. Nhập/chỉnh sửa thông tin (tên, ảnh đại diện,\ldots). & \\
\hline
5. Nhấn ``Lưu''. & \\
\hline
& 6. Kiểm tra dữ liệu hợp lệ. \\
\hline
& 7. Lưu thông tin vào database. \\
\hline
& 8. Hiển thị thông báo cập nhật thành công. \\
\hline
\multicolumn{2}{|p{0.97\textwidth}|}{\textbf{Luồng sự kiện phụ (Alternative Flows):}} \\
\hline
\multicolumn{2}{|p{0.97\textwidth}|}{6.1 Hiển thị thông báo lỗi (ví dụ: tên rỗng, ảnh sai định dạng,\ldots).} \\
\hline
\multicolumn{2}{|p{0.97\textwidth}|}{6.2 Quay lại bước 4.} \\
\hline
\caption{Mô tả chức năng UC03}
\end{xltabular}
\renewcommand{\arraystretch}{1}

\begin{figure}[H]
    \centering
    \begin{plantuml}
        @startuml
        !theme plain
        title Cập nhật hồ sơ cá nhân
        |Người dùng|
        start
        :1. Truy cập vào trang "Hồ sơ cá nhân";
        |Hệ thống|
        :2. Hiển thị thông tin hiện tại\n(tên, ảnh, email, ...);
        |Người dùng|
        :3. Chọn "Chỉnh sửa sơ";
        :4. Nhập/chỉnh sửa thông tin\n(tên, ảnh đại diện, ...);
        :5. Nhấn "Lưu";
        |Hệ thống|
        if (6. Kiểm tra dữ liệu hợp lệ) then (Hợp lệ)
          :7. Lưu thông tin vào database;
          :8. Hiển thị thông báo cập nhật thành công;
        else (Không hợp lệ)
          #pink:6.1 Hiển thị thông báo lỗi\n(ví dụ: tên rỗng, ảnh sai định dạng, ...);
          :6.2 Quay lại bước 4;
          detach
        endif
        stop
        @enduml
    \end{plantuml}
    \caption{Activity diagram --- Cập nhật hồ sơ cá nhân}
\end{figure}

\begin{figure}[H]
    \centering
    \begin{plantuml}
        @startuml
        !theme plain
        actor "Người dùng" as U
        participant "GD_HoSoCaNhan\n(Giao diện)" as GD
        participant "Ctrl_HoSoCaNhan\n(Controller)" as CTRL
        participant "UserService" as US
        participant "UserModel" as UM
        database "Database" as DB

        U -> GD: 1. Truy cập trang "Hồ sơ cá nhân"
        GD -> CTRL: 1.1 Yêu cầu lấy thông tin hồ sơ
        activate CTRL
        CTRL -> US: 1.1.1 getProfile(userId)
        activate US
        US -> UM: 1.1.1.1 findById(userId)
        activate UM
        UM -> DB: 1.1.1.1.1 SELECT * FROM users WHERE id = userId
        activate DB
        DB --> UM: 1.1.1.1.2 return user data
        deactivate DB
        UM --> US: 1.1.1.2 return user data
        deactivate UM
        US --> CTRL: 1.1.2 return profile data
        deactivate US
        CTRL --> GD: 1.2 Hiển thị thông tin hiện tại (tên, ảnh, email, ...)
        deactivate CTRL

        U -> GD: 2. Chọn "Chỉnh sửa hồ sơ"
        GD --> U: 2.1 Hiển thị form chỉnh sửa (đã điền sẵn thông tin)

        U -> GD: 3. Nhập/chỉnh sửa thông tin (tên, ảnh đại diện, ...)
        U -> GD: 4. Nhấn "Lưu"
        GD -> CTRL: 4.1 Gửi dữ liệu cập nhật (userId, data)
        activate CTRL
        CTRL -> US: 4.1.1 updateProfile(userId, data)
        activate US
        US -> UM: 4.1.1.1 update(userId, data)
        activate UM
        UM -> DB: 4.1.1.1.1 UPDATE users SET ... WHERE id = userId
        activate DB
        DB --> UM: 4.1.1.1.2 return result
        deactivate DB
        UM --> US: 4.1.1.2 return result
        deactivate UM
        US --> CTRL: 4.1.2 return result
        deactivate US

        alt [Thành công]
          CTRL --> GD: 5.1 Hiển thị thông báo thành công
          GD --> U: 5.1 Thông báo cập nhật thành công
        else [Thất bại]
          CTRL --> GD: 5.2 Hiển thị thông báo lỗi\n(ví dụ: tên rỗng, ảnh sai định dạng, ...)
          GD --> U: 5.4 Hiển thị lại form chỉnh sửa để người dùng sửa lại
        end
        deactivate CTRL
        @enduml
    \end{plantuml}
    \caption{Sequence diagram --- Cập nhật hồ sơ cá nhân}
\end{figure}

\subsubsection{UC04 --- Tìm kiếm người dùng}

\renewcommand{\arraystretch}{1.4}
\begin{xltabular}{\textwidth}{|X|X|}
\hline
\multicolumn{2}{|p{0.97\textwidth}|}{\textbf{Tên usecase:} Tìm kiếm người dùng.} \\
\hline
\multicolumn{2}{|p{0.97\textwidth}|}{\textbf{Actor:} Người dùng.} \\
\hline
\multicolumn{2}{|p{0.97\textwidth}|}{\textbf{Mô tả:} Cho phép người dùng tìm kiếm tài khoản khác bằng tên, email.} \\
\hline
\multicolumn{2}{|p{0.97\textwidth}|}{\textbf{Tiền điều kiện:} Người dùng đã đăng nhập vào hệ thống.} \\
\hline
\multicolumn{2}{|p{0.97\textwidth}|}{\textbf{Hậu điều kiện:}} \\
\hline
\multicolumn{2}{|p{0.97\textwidth}|}{\quad $+$ \textbf{Thành công:} Danh sách người dùng phù hợp được hiển thị.} \\
\hline
\multicolumn{2}{|p{0.97\textwidth}|}{\quad $+$ \textbf{Thất bại:} Hiển thị thông báo không tìm thấy hoặc lỗi hệ thống.} \\
\hline
\multicolumn{2}{|c|}{\textbf{Luồng sự kiện chính (Basic flows):}} \\
\hline
\multicolumn{1}{|c|}{\textbf{Actor}} & \multicolumn{1}{c|}{\textbf{System}} \\
\hline
1. Người dùng truy cập vào chức năng tìm kiếm. & \\
\hline
2. Nhập từ khóa tìm kiếm (tên/email). & \\
\hline
3. Nhấn ``Tìm kiếm''. & \\
\hline
& 4. Nhận yêu cầu tìm kiếm. \\
\hline
& 5. Tìm kiếm dữ liệu người dùng trong database. \\
\hline
& 6. Hiển thị danh sách người dùng. \\
\hline
\multicolumn{2}{|p{0.97\textwidth}|}{\textbf{Luồng sự kiện phụ (Alternative Flows):}} \\
\hline
\multicolumn{2}{|p{0.97\textwidth}|}{5.1 Không tìm thấy kết quả.} \\
\hline
\multicolumn{2}{|p{0.97\textwidth}|}{5.2 Hiển thị thông báo: ``Không tìm thấy người dùng''.} \\
\hline
\multicolumn{2}{|p{0.97\textwidth}|}{5.3 Cho phép nhập lại từ khóa (quay về bước 2).} \\
\hline
\caption{Mô tả chức năng UC04}
\end{xltabular}
\renewcommand{\arraystretch}{1}

\begin{figure}[H]
    \centering
    \begin{plantuml}
        @startuml
        !theme plain
        title Tìm kiếm người dùng
        |Người dùng|
        start
        :1. Người dùng truy cập\nvào chức năng tìm kiếm;
        :2. Nhập từ khóa tìm kiếm (tên/email);
        :3. Nhấn "Tìm kiếm";
        |Hệ thống|
        :4. Nhận yêu cầu tìm kiếm;
        :5. Tìm kiếm dữ liệu người dùng\ntrong database;
        if (5.1. Có tìm thấy kết quả?) then (Có)
          :6. Hiển thị danh sách người dùng phù hợp;
        else (Không)
          #pink:5.2 Hiển thị thông báo:\n"Không tìm thấy người dùng";
          :5.3 Cho phép nhập lại từ khóa\n(quay về bước 2);
          detach
        endif
        stop
        @enduml
    \end{plantuml}
    \caption{Activity diagram --- Tìm kiếm người dùng}
\end{figure}

\begin{figure}[H]
    \centering
    \begin{plantuml}
        @startuml
        !theme plain
        actor "Người dùng" as U
        participant "GD_TimKiem\n(Giao diện)" as GD
        participant "Ctrl_TimKiem\n(Controller)" as CTRL
        participant "UserService" as US
        participant "UserModel" as UM
        database "Database" as DB

        U -> GD: 1. Truy cập chức năng tìm kiếm
        U -> GD: 2. Nhập từ khóa (tên/email)
        U -> GD: 3. Nhấn "Tìm kiếm"
        GD -> CTRL: 4. Gửi request tìm kiếm (keyword)
        activate CTRL
        CTRL -> US: 5. search(keyword)
        activate US
        US -> UM: 5.1 findAll(keyword)
        activate UM
        UM -> DB: SELECT * FROM users\nWHERE name/email LIKE keyword
        activate DB
        DB --> UM: result
        deactivate DB
        UM --> US: danh sách user
        deactivate UM
        US --> CTRL: result
        deactivate US

        alt [Có kết quả]
          CTRL --> GD: 6. Trả danh sách người dùng phù hợp
          GD --> U: 6. Hiển thị danh sách người dùng
        else [Không có kết quả]
          CTRL --> GD: 5.2 Thông báo "Không tìm thấy"
          GD --> U: 5.2 Hiển thị thông báo
        end
        deactivate CTRL
        @enduml
    \end{plantuml}
    \caption{Sequence diagram --- Tìm kiếm người dùng}
\end{figure}

\subsubsection{UC05 --- Gửi lời mời kết bạn}

\renewcommand{\arraystretch}{1.4}
\begin{xltabular}{\textwidth}{|X|X|}
\hline
\multicolumn{2}{|p{0.97\textwidth}|}{\textbf{Tên usecase:} Gửi lời mời kết bạn.} \\
\hline
\multicolumn{2}{|p{0.97\textwidth}|}{\textbf{Actor:} Người dùng.} \\
\hline
\multicolumn{2}{|p{0.97\textwidth}|}{\textbf{Mô tả:} Cho phép người dùng gửi lời mời kết bạn đến một người dùng khác.} \\
\hline
\multicolumn{2}{|p{0.97\textwidth}|}{\textbf{Tiền điều kiện:} Người dùng đã đăng nhập, đã tìm kiếm và xem được profile của người dùng khác (UC04).} \\
\hline
\multicolumn{2}{|p{0.97\textwidth}|}{\textbf{Hậu điều kiện:}} \\
\hline
\multicolumn{2}{|p{0.97\textwidth}|}{\quad $+$ \textbf{Thành công:} Lời mời kết bạn được gửi, người nhận nhận được thông báo qua WebSocket.} \\
\hline
\multicolumn{2}{|p{0.97\textwidth}|}{\quad $+$ \textbf{Thất bại:} Hiển thị thông báo lỗi tương ứng (đã là bạn, đã gửi lời mời, bị chặn).} \\
\hline
\multicolumn{2}{|c|}{\textbf{Luồng sự kiện chính (Basic flows):}} \\
\hline
\multicolumn{1}{|c|}{\textbf{Actor}} & \multicolumn{1}{c|}{\textbf{System}} \\
\hline
1. Người dùng xem profile của người dùng khác. & \\
\hline
2. Nhấn nút ``Kết bạn''. & \\
\hline
& 3. Nhận yêu cầu POST /api/v1/friend-request với targetUserId. \\
\hline
& 4. Kiểm tra quan hệ bạn bè, lời mời đã tồn tại, và trạng thái chặn. \\
\hline
& 5. Nếu hợp lệ, tạo friend request với trạng thái PENDING. \\
\hline
& 6. Gửi WebSocket notification đến /user/\{email\}/queue/friend-request cho người nhận. \\
\hline
& 7. Trả về phản hồi thành công cho frontend. \\
\hline
8. Thấy thông báo ``Đã gửi lời mời kết bạn''. & \\
\hline
\multicolumn{2}{|p{0.97\textwidth}|}{\textbf{Luồng sự kiện phụ (Alternative Flows):}} \\
\hline
\multicolumn{2}{|p{0.97\textwidth}|}{4.1 Đã là bạn bè: Backend trả về lỗi 409 Conflict.} \\
\hline
\multicolumn{2}{|p{0.97\textwidth}|}{4.2 Frontend hiển thị ``Các bạn đã là bạn bè''.} \\
\hline
\multicolumn{2}{|p{0.97\textwidth}|}{4.3 Đã gửi lời mời trước đó: Backend trả về lỗi 409 Conflict.} \\
\hline
\multicolumn{2}{|p{0.97\textwidth}|}{4.4 Frontend hiển thị ``Đã gửi lời mời trước đó''.} \\
\hline
\multicolumn{2}{|p{0.97\textwidth}|}{4.5 Người dùng mục tiêu bị chặn: Backend từ chối tạo friend request.} \\
\hline
\caption{Mô tả chức năng UC05}
\end{xltabular}
\renewcommand{\arraystretch}{1}

\begin{figure}[H]
    \centering
    \begin{plantuml}
        @startuml
        !theme plain
        |User A|
        start
        :Xem profile người dùng B;
        :Nhấn "Kết bạn";
        |Frontend|
        :POST /friend-request;
        |Backend|
        :Kiểm tra quan hệ bạn bè;
        |Database|
        :SELECT friend_request, block;
        |Backend|
        if (Đã là bạn?) then (Có - ngoại lệ)
          |Frontend|
          :Hiển thị "Đã là bạn bè";
          |User A|
          stop
        else (Chưa là bạn)
          if (Đã gửi lời mời?) then (Rồi - ngoại lệ)
            |Frontend|
            :Hiển thị "Đã gửi lời mời";
            |User A|
            stop
          else (Chưa gửi)
            if (Bị chặn?) then (Bị chặn - ngoại lệ)
              |Frontend|
              :Hiển thị "Không thể gửi lời mời";
              |User A|
              stop
            else (OK)
              :INSERT friend_request (PENDING);
              |Database|
              :Lưu lời mời;
              |Backend|
              :Gửi WS notification;
              |WebSocket|
              :/user/queue/friend-request;
              |User B|
              :Nhận thông báo lời mời;
              |User A|
              :Thấy "Đã gửi lời mời";
              stop
            endif
          endif
        endif
        stop
        @enduml
    \end{plantuml}
    \caption{Activity diagram --- Gửi lời mời kết bạn}
\end{figure}

\begin{figure}[H]
    \centering
    \begin{plantuml}
        @startuml
        !theme plain
        actor "User A" as U
        participant "Frontend A" as FE
        participant "Backend" as BE
        database "MySQL" as DB
        participant "WebSocket" as WS
        actor "User B" as R

        U -> FE: Nhấn "Kết bạn" trên profile B
        activate FE
        FE -> BE: POST /friend-request \n{ targetUserId: B }
        activate BE

        BE -> DB: Kiểm tra bạn bè, lời mời, block
        activate DB
        DB --> BE: { relation status }
        deactivate DB

        alt Chưa có quan hệ
          BE -> DB: INSERT friend_request (PENDING)
          activate DB
          DB --> BE: Saved
          deactivate DB
          BE -> WS: SEND /user/{email}/queue/friend-request
          activate WS
          WS --> R: Thông báo lời mời kết bạn
          deactivate WS
          BE --> FE: { id, status: PENDING }
          deactivate BE
          FE --> U: "Đã gửi lời mời kết bạn"
        else Đã là bạn bè
          BE --> FE: 409 Conflict
          deactivate BE
          FE --> U: "Các bạn đã là bạn bè"
        else Đã gửi lời mời trước đó
          BE --> FE: 409 Conflict
          deactivate BE
          FE --> U: "Đã gửi lời mời trước đó"
        end
        deactivate FE
        @enduml
    \end{plantuml}
    \caption{Sequence diagram --- Gửi lời mời kết bạn}
\end{figure}

\subsubsection{UC06 --- Chấp nhận lời mời kết bạn}

\renewcommand{\arraystretch}{1.4}
\begin{xltabular}{\textwidth}{|X|X|}
\hline
\multicolumn{2}{|p{0.97\textwidth}|}{\textbf{Tên usecase:} Chấp nhận lời mời kết bạn.} \\
\hline
\multicolumn{2}{|p{0.97\textwidth}|}{\textbf{Actor:} Người dùng.} \\
\hline
\multicolumn{2}{|p{0.97\textwidth}|}{\textbf{Mô tả:} Cho phép người dùng chấp nhận lời mời kết bạn từ người dùng khác.} \\
\hline
\multicolumn{2}{|p{0.97\textwidth}|}{\textbf{Tiền điều kiện:} Người dùng đã đăng nhập, có lời mời kết bạn đang ở trạng thái PENDING từ người dùng khác.} \\
\hline
\multicolumn{2}{|p{0.97\textwidth}|}{\textbf{Hậu điều kiện:}} \\
\hline
\multicolumn{2}{|p{0.97\textwidth}|}{\quad $+$ \textbf{Thành công:} Hai người trở thành bạn bè, chat 1-1 được tạo, cả hai nhận được thông báo.} \\
\hline
\multicolumn{2}{|p{0.97\textwidth}|}{\quad $+$ \textbf{Thất bại:} Hiển thị thông báo lỗi (lời mời không tồn tại, đã xử lý, hoặc không có quyền).} \\
\hline
\multicolumn{2}{|c|}{\textbf{Luồng sự kiện chính (Basic flows):}} \\
\hline
\multicolumn{1}{|c|}{\textbf{Actor}} & \multicolumn{1}{c|}{\textbf{System}} \\
\hline
1. Người dùng nhận thông báo lời mời kết bạn qua WebSocket. & \\
\hline
2. Mở danh sách lời mời kết bạn. & \\
\hline
& 3. Nhận yêu cầu GET /api/v1/friend-request/pending. \\
\hline
& 4. Truy vấn và trả về danh sách lời mời đang ở trạng thái PENDING. \\
\hline
5. Chọn lời mời và nhấn ``Chấp nhận''. & \\
\hline
& 6. Nhận yêu cầu PUT /api/v1/friend-request/\{id\}/accept. \\
\hline
& 7. Cập nhật trạng thái friend request thành ACCEPTED. \\
\hline
& 8. Tự động tạo chat 1-1 giữa hai người dùng. \\
\hline
& 9. Gửi WebSocket notification đến người gửi lời mời. \\
\hline
& 10. Trả về phản hồi thành công cho frontend. \\
\hline
11. Thấy thông báo ``Đã kết bạn thành công''. & \\
\hline
\multicolumn{2}{|p{0.97\textwidth}|}{\textbf{Luồng sự kiện phụ (Alternative Flows):}} \\
\hline
\multicolumn{2}{|p{0.97\textwidth}|}{6.1 Lời mời không tồn tại: Backend trả về lỗi 404 Not Found.} \\
\hline
\multicolumn{2}{|p{0.97\textwidth}|}{6.2 Lời mời đã được xử lý (đã chấp nhận/từ chối trước đó): Backend trả về lỗi 400 Bad Request.} \\
\hline
\multicolumn{2}{|p{0.97\textwidth}|}{6.3 Người dùng không phải là người nhận: Backend trả về lỗi 403 Forbidden.} \\
\hline
\caption{Mô tả chức năng UC06}
\end{xltabular}
\renewcommand{\arraystretch}{1}

\begin{figure}[H]
    \centering
    \begin{plantuml}
        @startuml
        !theme plain
        |User B|
        start
        :Nhận thông báo lời mời;
        :Mở danh sách lời mời;
        |Frontend|
        :GET /friend-request/pending;
        |Backend|
        :Lấy danh sách PENDING;
        |Database|
        :SELECT pending requests;
        |Backend|
        |Frontend|
        :Hiển thị danh sách;
        |User B|
        :Chọn lời mời từ A;
        :Nhấn "Chấp nhận";
        |Frontend|
        :PUT /friend-request/{id}/accept;
        |Backend|
        if (Lời mời tồn tại?) then (Có)
          if (Trạng thái PENDING?) then (Còn hiệu lực)
            :UPDATE status = ACCEPTED;
            :INSERT chat 1-1;
            |Database|
            :Cập nhật DB;
            |Backend|
            :Gửi WS notification;
            |WebSocket|
            :/user/queue/friend-request-accepted;
            |User A|
            :Nhận thông báo;
            |User B|
            :Thấy "Đã kết bạn";
            stop
          else (Đã xử lý - ngoại lệ)
            |Frontend|
            :Hiển thị "Lời mời đã xử lý";
            |User B|
            stop
          endif
        else (Không - ngoại lệ)
          |Frontend|
          :Hiển thị "Không tìm thấy";
          |User B|
          stop
        endif
        stop
        @enduml
    \end{plantuml}
    \caption{Activity diagram --- Chấp nhận lời mời kết bạn}
\end{figure}

\begin{figure}[H]
    \centering
    \begin{plantuml}
        @startuml
        !theme plain
        actor "User B (Người nhận)" as U
        participant "Frontend B" as FE
        participant "Backend" as BE
        database "MySQL" as DB
        participant "WebSocket" as WS
        actor "User A (Người gửi)" as S

        U -> FE: Mở danh sách lời mời
        activate FE
        FE -> BE: GET /friend-request/pending
        activate BE
        BE -> DB: SELECT * FROM friend_request \nWHERE status = PENDING
        activate DB
        DB --> BE: [{ id, from: A }]
        deactivate DB
        BE --> FE: Danh sách lời mời
        deactivate BE
        FE --> U: Hiển thị lời mời
        deactivate FE

        U -> FE: Nhấn "Chấp nhận" lời mời từ A
        activate FE
        FE -> BE: PUT /friend-request/{id}/accept
        activate BE
        BE -> DB: SELECT request by id
        activate DB

        alt Request tồn tại và PENDING
          DB --> BE: { status: PENDING }
          deactivate DB
          BE -> DB: UPDATE status = ACCEPTED
          activate DB
          DB --> BE: OK
          deactivate DB
          BE -> DB: INSERT chat \n(userA, userB)
          activate DB
          DB --> BE: chat created
          deactivate DB
          BE -> WS: SEND friend-request-accepted
          activate WS
          WS --> S: Thông báo "B đã chấp nhận"
          deactivate WS
          BE --> FE: { status: ACCEPTED, chatId }
          deactivate BE
          FE --> U: "Đã kết bạn"
        else Request không tồn tại
          DB --> BE: null
          deactivate DB
          BE --> FE: 404 Not Found
          deactivate BE
          FE --> U: "Lời mời không tồn tại"
        else Request đã xử lý
          DB --> BE: { status: REJECTED }
          deactivate DB
          BE --> FE: 400 Bad Request
          deactivate BE
          FE --> U: "Lời mời đã được xử lý"
        end
        deactivate FE
        @enduml
    \end{plantuml}
    \caption{Sequence diagram --- Chấp nhận lời mời kết bạn}
\end{figure}

\subsubsection{UC07 --- Nhắn tin văn bản 1-1}

\renewcommand{\arraystretch}{1.4}
\begin{xltabular}{\textwidth}{|X|X|}
\hline
\multicolumn{2}{|p{0.97\textwidth}|}{\textbf{Tên usecase:} Nhắn tin văn bản 1-1.} \\
\hline
\multicolumn{2}{|p{0.97\textwidth}|}{\textbf{Actor:} Người dùng.} \\
\hline
\multicolumn{2}{|p{0.97\textwidth}|}{\textbf{Mô tả:} Cho phép người dùng gửi tin nhắn văn bản đến một người dùng khác trong cuộc hội thoại 1-1.} \\
\hline
\multicolumn{2}{|p{0.97\textwidth}|}{\textbf{Tiền điều kiện:} Người dùng đã đăng nhập, có quan hệ bạn bè với người nhận và có chat 1-1 tồn tại.} \\
\hline
\multicolumn{2}{|p{0.97\textwidth}|}{\textbf{Hậu điều kiện:}} \\
\hline
\multicolumn{2}{|p{0.97\textwidth}|}{\quad $+$ \textbf{Thành công:} Tin nhắn được lưu và broadcast đến cả hai người dùng qua WebSocket.} \\
\hline
\multicolumn{2}{|p{0.97\textwidth}|}{\quad $+$ \textbf{Thất bại:} Hiển thị thông báo lỗi (nội dung trống, bị chặn, chat không tồn tại).} \\
\hline
\multicolumn{2}{|c|}{\textbf{Luồng sự kiện chính (Basic flows):}} \\
\hline
\multicolumn{1}{|c|}{\textbf{Actor}} & \multicolumn{1}{c|}{\textbf{System}} \\
\hline
1. Người dùng mở chat với người dùng khác từ danh sách chat. & \\
\hline
& 2. Nhận yêu cầu GET /api/v1/message/chat/\{chatId\}. \\
\hline
& 3. Truy vấn và trả về lịch sử tin nhắn. \\
\hline
4. Nhập nội dung tin nhắn và nhấn ``Gửi''. & \\
\hline
& 5. Nhận yêu cầu POST /api/v1/message với chatId và nội dung. \\
\hline
& 6. Kiểm tra quyền gửi (không bị chặn, chat tồn tại). \\
\hline
& 7. Lưu tin nhắn vào database với trạng thái SENT. \\
\hline
& 8. Broadcast tin nhắn đến /topic/chat/\{chatId\}. \\
\hline
& 9. Gửi notification đến /user/\{email\}/queue/messages của người nhận. \\
\hline
10. Nhận tin nhắn real-time qua WebSocket. & \\
\hline
\multicolumn{2}{|p{0.97\textwidth}|}{\textbf{Luồng sự kiện phụ (Alternative Flows):}} \\
\hline
\multicolumn{2}{|p{0.97\textwidth}|}{5.1 Nội dung tin nhắn trống: Backend trả về lỗi 400 Bad Request.} \\
\hline
\multicolumn{2}{|p{0.97\textwidth}|}{6.1 Người gửi bị người nhận chặn: Backend từ chối gửi, trả về lỗi 403 Forbidden.} \\
\hline
\multicolumn{2}{|p{0.97\textwidth}|}{6.2 Token hết hạn: Frontend tự động refresh token và thử lại.} \\
\hline
\multicolumn{2}{|p{0.97\textwidth}|}{6.3 Chat không tồn tại: Backend trả về lỗi 404 Not Found.} \\
\hline
\caption{Mô tả chức năng UC07}
\end{xltabular}
\renewcommand{\arraystretch}{1}

\begin{figure}[H]
    \centering
    \begin{plantuml}
        @startuml
        !theme plain
        |User A|
        start
        :Mở chat với B;
        |Frontend|
        :GET /message/chat/{chatId};
        :Hiển thị lịch sử tin nhắn;
        |User A|
        :Nhập nội dung;
        :Nhấn Gửi;
        |Frontend|
        if (Nội dung trống?) then (Có - ngoại lệ)
          :Hiển thị "Không thể gửi tin nhắn trống";
          |User A|
          stop
        else (Có nội dung)
          :POST /api/v1/message;
          |Backend|
          :Kiểm tra quyền gửi;
          if (Bị chặn?) then (Có - ngoại lệ)
            |Frontend|
            :Hiển thị "Bạn đã bị chặn";
            |User A|
            stop
          else (Không bị chặn)
            if (Chat tồn tại?) then (OK)
            :INSERT message (SENT);
            |Database|
            :Lưu tin nhắn;
            |Backend|
            :Gửi WS broadcast;
            :Gửi WS notification;
            |WebSocket|
            :/topic/chat/{id};
            :/user/queue/messages;
            |User A|
            :Tin nhắn hiển thị;
            |User B|
            :Nhận tin nhắn real-time;
            stop
            else (Không - ngoại lệ)
              |Frontend|
              :Hiển thị "Chat không tồn tại";
              |User A|
              stop
          endif
        endif
        endif
        stop
        @enduml
    \end{plantuml}
    \caption{Activity diagram --- Nhắn tin văn bản 1-1}
\end{figure}

\begin{figure}[H]
    \centering
    \begin{plantuml}
        @startuml
        !theme plain
        actor "Người gửi (Sender)" as S
        participant "Frontend" as FE
        participant "Backend" as BE
        database "Database" as DB
        participant "WebSocket Broker" as WS
        actor "Người nhận (Receiver)" as R

        S -> FE: Nhập nội dung + nhấn Gửi
        activate FE
        FE -> BE: POST /api/v1/message \n{ chatId, content }
        activate BE

        BE -> BE: Kiểm tra quyền gửi\n(không bị chặn, chat tồn tại)

        alt Gửi hợp lệ
          BE -> DB: INSERT message (SENT)
          activate DB
          DB --> BE: message saved
          deactivate DB

          BE -> WS: SEND /topic/chat/{id}
          activate WS
          WS --> FE: MESSAGE (broadcast)
          WS --> R: MESSAGE (notification)
          deactivate WS

          BE -> WS: SEND /user/{email}/queue/messages
          BE --> FE: { id, content, timestamp }
          deactivate BE
          FE --> S: Hiển thị tin nhắn trong chat
          deactivate FE
          R --> R: Cập nhật danh sách chat
        else Bị chặn
          BE --> FE: 403 Blocked
          deactivate BE
          FE --> S: "Bạn đã bị chặn"
          deactivate FE
        else Chat không tồn tại
          BE --> FE: 404 Not Found
          deactivate BE
          FE --> S: "Cuộc trò chuyện không tồn tại"
          deactivate FE
        end
        @enduml
    \end{plantuml}
    \caption{Sequence diagram --- Nhắn tin văn bản 1-1}
\end{figure}

\subsubsection{UC08 --- Gửi file đa phương tiện}

\renewcommand{\arraystretch}{1.4}
\begin{xltabular}{\textwidth}{|X|X|}
\hline
\multicolumn{2}{|p{0.97\textwidth}|}{\textbf{Tên usecase:} Gửi file đa phương tiện.} \\
\hline
\multicolumn{2}{|p{0.97\textwidth}|}{\textbf{Actor:} Người dùng.} \\
\hline
\multicolumn{2}{|p{0.97\textwidth}|}{\textbf{Mô tả:} Cho phép người dùng gửi hình ảnh, video, âm thanh hoặc file đính kèm trong cuộc trò chuyện.} \\
\hline
\multicolumn{2}{|p{0.97\textwidth}|}{\textbf{Tiền điều kiện:} Người dùng đã đăng nhập, đang trong một cuộc trò chuyện (1-1 hoặc nhóm).} \\
\hline
\multicolumn{2}{|p{0.97\textwidth}|}{\textbf{Hậu điều kiện:}} \\
\hline
\multicolumn{2}{|p{0.97\textwidth}|}{\quad $+$ \textbf{Thành công:} File được upload lên Cloudflare R2, tin nhắn media được broadcast đến tất cả thành viên.} \\
\hline
\multicolumn{2}{|p{0.97\textwidth}|}{\quad $+$ \textbf{Thất bại:} Hiển thị thông báo lỗi (định dạng không hỗ trợ, file quá lớn, upload thất bại).} \\
\hline
\multicolumn{2}{|c|}{\textbf{Luồng sự kiện chính (Basic flows):}} \\
\hline
\multicolumn{1}{|c|}{\textbf{Actor}} & \multicolumn{1}{c|}{\textbf{System}} \\
\hline
1. Người dùng mở chat và chọn biểu tượng đính kèm. & \\
\hline
2. Chọn file từ máy tính (có thể chọn nhiều file). & \\
\hline
& 3. Kiểm tra loại file (ảnh, video, audio, PDF, Office, zip, txt) và kích thước (tối đa 50MB). \\
\hline
& 4. Nhận file qua POST /api/v1/message/upload-media/\{chatId\} (multipart/form-data). \\
\hline
& 5. Tạo UUID key cho file. \\
\hline
& 6. Upload file lên Cloudflare R2 (S3-compatible) với UUID key. \\
\hline
& 7. Lưu message với loại IMAGE/VIDEO/AUDIO/FILE và media URL. \\
\hline
& 8. Trả về message object chứa media URL cho frontend. \\
\hline
& 9. Broadcast message qua /topic/chat/\{chatId\}. \\
\hline
10. Xem media hiển thị trong luồng chat. & \\
\hline
\multicolumn{2}{|p{0.97\textwidth}|}{\textbf{Luồng sự kiện phụ (Alternative Flows):}} \\
\hline
\multicolumn{2}{|p{0.97\textwidth}|}{3.1 Định dạng file không được hỗ trợ: Frontend hiển thị thông báo lỗi, không cho phép gửi.} \\
\hline
\multicolumn{2}{|p{0.97\textwidth}|}{3.2 File vượt quá 50MB: Backend trả về lỗi 413 Payload Too Large.} \\
\hline
\multicolumn{2}{|p{0.97\textwidth}|}{6.1 Upload lên R2 thất bại: Backend trả về lỗi 500 Internal Server Error, frontend cho phép thử lại.} \\
\hline
\caption{Mô tả chức năng UC08}
\end{xltabular}
\renewcommand{\arraystretch}{1}

\begin{figure}[H]
    \centering
    \begin{plantuml}
        @startuml
        !theme plain
        |User|
        start
        :Mở chat;
        :Chọn biểu tượng đính kèm;
        :Chọn file từ máy tính;
        |Frontend|
        if (Loại file hợp lệ?) then (Hợp lệ)
          if (Kích thước <= 50MB?) then (OK)
            :POST /upload-media/{chatId};
            |Backend|
            :Tạo UUID key;
            :Upload file lên R2;
            |Cloudflare R2|
            :PUT object (S3);
            if (Upload thành công?) then (Có)
              |Backend|
              :INSERT message (IMAGE/VIDEO/FILE);
              |Database|
              :Lưu message + media URL;
              |Backend|
              :Gửi WS broadcast;
              |Frontend|
              :Hiển thị media trong chat;
              |User|
              stop
            else (Không - ngoại lệ)
              |Backend|
              |Frontend|
              :Hiển thị "Upload thất bại\nThử lại";
              |User|
              stop
            endif
          else (Quá lớn - ngoại lệ)
            :Hiển thị "File quá lớn";
            |User|
            stop
          endif
        else (Không - ngoại lệ)
          :Hiển thị "Định dạng không hỗ trợ";
          |User|
          stop
        endif
        stop
        @enduml
    \end{plantuml}
    \caption{Activity diagram --- Gửi file đa phương tiện}
\end{figure}

\begin{figure}[H]
    \centering
    \begin{plantuml}
        @startuml
        !theme plain
        actor "User" as U
        participant "Frontend" as FE
        participant "Backend" as BE
        participant "Cloudflare R2" as R2
        database "MySQL" as DB

        U -> FE: Chọn file từ máy
        activate FE
        FE -> FE: Kiểm tra loại file\nvà kích thước

        alt File hợp lệ
          FE -> BE: POST /api/v1/message/upload-media/{chatId}
          activate BE
          BE -> BE: Tạo UUID key cho file
          BE -> R2: PUT object (S3 upload)
          activate R2

          alt Upload thành công
            R2 --> BE: Upload success
            deactivate R2
            BE -> DB: INSERT message\n(type, mediaUrl, fileName)
            activate DB
            DB --> BE: message saved
            deactivate DB
            BE --> FE: { messageId, mediaUrl }
            deactivate BE
            FE --> U: Hiển thị media trong chat
          else Upload thất bại
            R2 --> BE: Error
            deactivate R2
            BE --> FE: 500 Upload failed
            deactivate BE
            FE --> U: "Tải file thất bại"
          end
        else File không hợp lệ
          FE --> U: "File không được hỗ trợ\nhoặc quá lớn"
        end
        deactivate FE
        @enduml
    \end{plantuml}
    \caption{Sequence diagram --- Gửi file đa phương tiện}
\end{figure}

\subsubsection{UC09 --- Tạo nhóm chat}

\renewcommand{\arraystretch}{1.4}
\begin{xltabular}{\textwidth}{|X|X|}
\hline
\multicolumn{2}{|p{0.97\textwidth}|}{\textbf{Tên usecase:} Tạo nhóm chat.} \\
\hline
\multicolumn{2}{|p{0.97\textwidth}|}{\textbf{Actor:} Người dùng.} \\
\hline
\multicolumn{2}{|p{0.97\textwidth}|}{\textbf{Mô tả:} Cho phép người dùng tạo một nhóm chat mới với nhiều thành viên.} \\
\hline
\multicolumn{2}{|p{0.97\textwidth}|}{\textbf{Tiền điều kiện:} Người dùng đã đăng nhập, có ít nhất một người bạn trong danh sách bạn bè.} \\
\hline
\multicolumn{2}{|p{0.97\textwidth}|}{\textbf{Hậu điều kiện:}} \\
\hline
\multicolumn{2}{|p{0.97\textwidth}|}{\quad $+$ \textbf{Thành công:} Nhóm chat được tạo, người tạo là ADMIN, các thành viên nhận được thông báo.} \\
\hline
\multicolumn{2}{|p{0.97\textwidth}|}{\quad $+$ \textbf{Thất bại:} Hiển thị thông báo lỗi (tên trống, không có thành viên, thành viên không hợp lệ).} \\
\hline
\multicolumn{2}{|c|}{\textbf{Luồng sự kiện chính (Basic flows):}} \\
\hline
\multicolumn{1}{|c|}{\textbf{Actor}} & \multicolumn{1}{c|}{\textbf{System}} \\
\hline
1. Người dùng chọn ``Tạo nhóm'' từ ứng dụng. & \\
\hline
& 2. Hiển thị form tạo nhóm (tên nhóm, danh sách bạn bè để chọn thành viên). \\
\hline
3. Nhập tên nhóm và chọn các thành viên từ danh sách bạn bè. & \\
\hline
& 4. Nhận yêu cầu POST /api/v1/group với tên nhóm và danh sách thành viên. \\
\hline
& 5. Tạo group và thêm người tạo với vai trò ADMIN, các thành viên với vai trò MEMBER. \\
\hline
& 6. Tự động gửi tin nhắn hệ thống ``[Tên] đã tạo nhóm''. \\
\hline
& 7. Broadcast group event đến các thành viên qua /topic/group/\{id\}/events. \\
\hline
8. Nhận thông báo và thấy nhóm mới trong danh sách. & \\
\hline
\multicolumn{2}{|p{0.97\textwidth}|}{\textbf{Luồng sự kiện phụ (Alternative Flows):}} \\
\hline
\multicolumn{2}{|p{0.97\textwidth}|}{4.1 Tên nhóm trống: Backend trả về lỗi 400 Bad Request.} \\
\hline
\multicolumn{2}{|p{0.97\textwidth}|}{4.2 Không chọn thành viên nào: Backend yêu cầu tối thiểu 1 thành viên ngoài người tạo.} \\
\hline
\multicolumn{2}{|p{0.97\textwidth}|}{4.3 Thành viên không phải bạn bè: Backend kiểm tra quan hệ bạn bè, bỏ qua hoặc báo lỗi.} \\
\hline
\caption{Mô tả chức năng UC09}
\end{xltabular}
\renewcommand{\arraystretch}{1}

\begin{figure}[H]
    \centering
    \begin{plantuml}
        @startuml
        !theme plain
        |User|
        start
        :Chọn "Tạo nhóm";
        |Frontend|
        :Hiển thị form\ntên nhóm + chọn thành viên;
        |User|
        :Nhập tên nhóm;
        :Chọn bạn bè làm thành viên;
        :Xác nhận tạo;
        |Frontend|
        if (Tên nhóm trống?) then (Có - ngoại lệ)
          :Hiển thị "Tên nhóm không được trống";
          |User|
          stop
        else (Có tên)
          if (Chưa chọn thành viên?) then (Có - ngoại lệ)
            :Hiển thị "Chọn ít nhất 1 thành viên";
            |User|
            stop
          else (OK)
            :POST /api/v1/group;
            |Backend|
            :Tạo group;
            :Thêm creator (ADMIN);
            :Thêm members (MEMBER);
            :Gửi tin nhắn hệ thống;
            |Database|
            :INSERT group, group_member;
            |Backend|
            :Gửi WS event;
            |WebSocket|
            :/topic/group/{id}/events;
            |Các thành viên|
            :Nhận thông báo;
            |User|
            :Vào màn hình nhóm mới;
            stop
          endif
        endif
        stop
        @enduml
    \end{plantuml}
    \caption{Activity diagram --- Tạo nhóm chat}
\end{figure}

\begin{figure}[H]
    \centering
    \begin{plantuml}
        @startuml
        !theme plain
        actor "User (Creator)" as C
        participant "Frontend" as FE
        participant "Backend" as BE
        database "MySQL" as DB
        participant "WebSocket" as WS
        actor "Các thành viên" as M

        C -> FE: Chọn "Tạo nhóm"
        activate FE
        C -> FE: Nhập tên + chọn thành viên
        FE -> BE: POST /api/v1/group \n{ name, memberIds }
        activate BE

        BE -> DB: INSERT group
        activate DB
        DB --> BE: group created
        deactivate DB
        BE -> DB: INSERT group_member\n(creator = ADMIN)
        activate DB
        DB --> BE: OK
        deactivate DB
        BE -> DB: INSERT group_member\n(members = MEMBER)
        activate DB
        DB --> BE: OK
        deactivate DB
        BE -> DB: INSERT system message\n"[User] đã tạo nhóm"
        activate DB
        DB --> BE: OK
        deactivate DB

        BE -> WS: SEND /topic/group/{id}/events\n{ type: GROUP_CREATED }
        activate WS
        WS --> M: Nhận thông báo nhóm mới
        deactivate WS
        BE --> FE: { id, name, members }
        deactivate BE
        FE --> C: Vào màn hình nhóm mới
        deactivate FE
        @enduml
    \end{plantuml}
    \caption{Sequence diagram --- Tạo nhóm chat}
\end{figure}

\subsubsection{UC10 --- Thêm thành viên vào nhóm}

\renewcommand{\arraystretch}{1.4}
\begin{xltabular}{\textwidth}{|X|X|}
\hline
\multicolumn{2}{|p{0.97\textwidth}|}{\textbf{Tên usecase:} Thêm thành viên vào nhóm.} \\
\hline
\multicolumn{2}{|p{0.97\textwidth}|}{\textbf{Actor:} Người dùng (Admin nhóm).} \\
\hline
\multicolumn{2}{|p{0.97\textwidth}|}{\textbf{Mô tả:} Cho phép thành viên có quyền (Admin nhóm) thêm người dùng khác vào nhóm chat.} \\
\hline
\multicolumn{2}{|p{0.97\textwidth}|}{\textbf{Tiền điều kiện:} Người dùng đã đăng nhập, là thành viên của nhóm với vai trò ADMIN hoặc có quyền mời.} \\
\hline
\multicolumn{2}{|p{0.97\textwidth}|}{\textbf{Hậu điều kiện:}} \\
\hline
\multicolumn{2}{|p{0.97\textwidth}|}{\quad $+$ \textbf{Thành công:} Thành viên mới được thêm vào nhóm, tất cả thành viên nhận được thông báo.} \\
\hline
\multicolumn{2}{|p{0.97\textwidth}|}{\quad $+$ \textbf{Thất bại:} Hiển thị thông báo lỗi (không có quyền, nhóm không tồn tại, thành viên đã có trong nhóm).} \\
\hline
\multicolumn{2}{|c|}{\textbf{Luồng sự kiện chính (Basic flows):}} \\
\hline
\multicolumn{1}{|c|}{\textbf{Actor}} & \multicolumn{1}{c|}{\textbf{System}} \\
\hline
1. Người dùng mở thông tin nhóm và chọn ``Thêm thành viên''. & \\
\hline
& 2. Hiển thị danh sách bạn bè chưa có trong nhóm. \\
\hline
3. Chọn một hoặc nhiều bạn bè và xác nhận. & \\
\hline
& 4. Nhận yêu cầu POST /api/v1/group/\{id\}/members với danh sách userId. \\
\hline
& 5. Kiểm tra quyền (người thêm phải là ADMIN hoặc có quyền mời). \\
\hline
& 6. Thêm các thành viên mới vào group\_member. \\
\hline
& 7. Gửi WebSocket event MEMBER\_ADDED đến /topic/group/\{id\}/events. \\
\hline
8. Nhận thông báo cập nhật danh sách thành viên. & \\
\hline
\multicolumn{2}{|p{0.97\textwidth}|}{\textbf{Luồng sự kiện phụ (Alternative Flows):}} \\
\hline
\multicolumn{2}{|p{0.97\textwidth}|}{5.1 Người dùng không có quyền thêm thành viên: Backend trả về lỗi 403 Forbidden.} \\
\hline
\multicolumn{2}{|p{0.97\textwidth}|}{5.2 Nhóm không tồn tại: Backend trả về lỗi 404 Not Found.} \\
\hline
\multicolumn{2}{|p{0.97\textwidth}|}{6.1 Người dùng đã có trong nhóm: Backend bỏ qua, không thêm trùng.} \\
\hline
\end{xltabular}
\renewcommand{\arraystretch}{1}

\begin{figure}[H]
    \centering
    \begin{plantuml}
        @startuml
        !theme plain
        |User (Admin)|
        start
        :Mở thông tin nhóm;
        :Chọn "Thêm thành viên";
        |Frontend|
        :GET danh sách bạn bè\nchưa trong nhóm;
        :Hiển thị danh sách;
        |User (Admin)|
        :Chọn bạn bè để thêm;
        :Xác nhận;
        |Frontend|
        :POST /group/{id}/members;
        |Backend|
        if (Là Admin nhóm?) then (Có)
          if (Nhóm tồn tại?) then (Có)
            :Thêm thành viên mới;
            |Database|
            :INSERT group_member;
            |Backend|
            :Gửi WS event;
            |WebSocket|
            :/topic/group/{id}/events\nMEMBER_ADDED;
            |Các thành viên|
            :Nhận cập nhật;
            |User (Admin)|
            :Thấy thông báo thành công;
            stop
          else (Không - ngoại lệ)
            |Frontend|
            :Hiển thị "Nhóm không tồn tại";
            |User (Admin)|
            stop
          endif
        else (Không - ngoại lệ)
          |Frontend|
          :Hiển thị "Không có quyền";
          |User (Admin)|
          stop
        endif
        stop
        @enduml
    \end{plantuml}
    \caption{Activity diagram --- Thêm thành viên vào nhóm}
\end{figure}

\begin{figure}[H]
    \centering
    \begin{plantuml}
        @startuml
        !theme plain
        actor "User (Admin)" as U
        participant "Frontend" as FE
        participant "Backend" as BE
        database "MySQL" as DB
        participant "WebSocket" as WS
        actor "Các thành viên" as M

        U -> FE: Chọn "Thêm thành viên"
        activate FE
        U -> FE: Chọn bạn + xác nhận
        FE -> BE: POST /group/{id}/members \n{ userIds }
        activate BE

        BE -> BE: Kiểm tra quyền ADMIN\n(trong group_member)

        alt Có quyền
          BE -> DB: INSERT group_member \n(userId, role=MEMBER)
          activate DB
          DB --> BE: Added
          deactivate DB
          BE -> WS: SEND /topic/group/{id}/events\n{ type: MEMBER_ADDED, userIds }
          activate WS
          WS --> M: Nhận thông báo\nthành viên mới
          deactivate WS
          BE --> FE: { success, added }
          deactivate BE
          FE --> U: "Đã thêm thành viên"
        else Không có quyền
          BE --> FE: 403 Forbidden
          deactivate BE
          FE --> U: "Bạn không có quyền thêm\nthành viên"
        end
        deactivate FE
        @enduml
    \end{plantuml}
    \caption{Sequence diagram --- Thêm thành viên vào nhóm}
\end{figure}

\subsubsection{UC11 --- Thả cảm xúc cho tin nhắn}

\renewcommand{\arraystretch}{1.4}
\begin{xltabular}{\textwidth}{|X|X|}
\hline
\multicolumn{2}{|p{0.97\textwidth}|}{\textbf{Tên usecase:} Thả cảm xúc cho tin nhắn.} \\
\hline
\multicolumn{2}{|p{0.97\textwidth}|}{\textbf{Actor:} Người dùng.} \\
\hline
\multicolumn{2}{|p{0.97\textwidth}|}{\textbf{Mô tả:} Cho phép người dùng thả cảm xúc (emoji reaction) cho một tin nhắn trong chat 1-1 hoặc nhóm.} \\
\hline
\multicolumn{2}{|p{0.97\textwidth}|}{\textbf{Tiền điều kiện:} Người dùng đã đăng nhập, đang xem tin nhắn trong một cuộc trò chuyện (1-1 hoặc nhóm).} \\
\hline
\multicolumn{2}{|p{0.97\textwidth}|}{\textbf{Hậu điều kiện:}} \\
\hline
\multicolumn{2}{|p{0.97\textwidth}|}{\quad $+$ \textbf{Thành công:} Reaction được thêm (hoặc gỡ nếu toggle off), broadcast đến tất cả người dùng trong chat.} \\
\hline
\multicolumn{2}{|p{0.97\textwidth}|}{\quad $+$ \textbf{Thất bại:} Hiển thị thông báo lỗi (tin nhắn không tồn tại, đã bị thu hồi).} \\
\hline
\multicolumn{2}{|c|}{\textbf{Luồng sự kiện chính (Basic flows):}} \\
\hline
\multicolumn{1}{|c|}{\textbf{Actor}} & \multicolumn{1}{c|}{\textbf{System}} \\
\hline
1. Người dùng di chuột/nhấn giữ vào một tin nhắn để hiện menu reaction. & \\
\hline
2. Chọn một emoji từ bảng chọn. & \\
\hline
& 3. Nhận yêu cầu POST /api/v1/reaction/message/\{messageId\} với emoji. \\
\hline
& 4. Kiểm tra tin nhắn tồn tại và chưa bị thu hồi/xóa. \\
\hline
& 5. Nếu reaction đã tồn tại (cùng user + cùng emoji): xóa reaction (toggle off). \\
\hline
& 6. Nếu chưa tồn tại: thêm reaction mới. \\
\hline
& 7. Broadcast cập nhật reactions qua WebSocket đến topic tương ứng. \\
\hline
8. Xem reactions được cập nhật trên tin nhắn. & \\
\hline
\multicolumn{2}{|p{0.97\textwidth}|}{\textbf{Luồng sự kiện phụ (Alternative Flows):}} \\
\hline
\multicolumn{2}{|p{0.97\textwidth}|}{4.1 Tin nhắn không tồn tại: Backend trả về lỗi 404 Not Found.} \\
\hline
\multicolumn{2}{|p{0.97\textwidth}|}{4.2 Tin nhắn đã bị thu hồi hoặc xóa: Backend từ chối thêm reaction.} \\
\hline
\multicolumn{2}{|p{0.97\textwidth}|}{4.3 Reaction cho group message: Backend xử lý tương tự qua POST /api/v1/reaction/group-message/\{id\}.} \\
\hline
\end{xltabular}
\renewcommand{\arraystretch}{1}

\begin{figure}[H]
    \centering
    \begin{plantuml}
        @startuml
        !theme plain
        |User|
        start
        :Chọn tin nhắn;
        :Mở menu reaction;
        :Chọn emoji;
        |Frontend|
        :POST /reaction/message/{id};
        |Backend|
        if (Tin nhắn tồn tại?) then (Có)
          if (Đã có reaction này?) then (Có - toggle off)
            :DELETE reaction;
            |Database|
            :Xóa reaction;
            |Backend|
          else (Chưa - thêm mới)
            :INSERT reaction;
            |Database|
            :Thêm reaction;
            |Backend|
          endif
          :Gửi WS broadcast reactions;
          |WebSocket|
          :/topic/chat/{id};
          |Frontend|
          :Cập nhật UI reactions;
          |User|
          :Thấy emoji trên tin nhắn;
          stop
        else (Không - ngoại lệ)
          |Frontend|
          :Hiển thị "Tin nhắn không tồn tại";
          |User|
          stop
        endif
        stop
        @enduml
    \end{plantuml}
    \caption{Activity diagram --- Thả cảm xúc cho tin nhắn}
\end{figure}

\begin{figure}[H]
    \centering
    \begin{plantuml}
        @startuml
        !theme plain
        actor "User" as U
        participant "Frontend" as FE
        participant "Backend" as BE
        database "MySQL" as DB
        participant "WebSocket" as WS

        U -> FE: Chọn emoji reaction
        activate FE
        FE -> BE: POST /reaction/message/{id}\n{ emoji: "heart" }
        activate BE

        BE -> DB: SELECT message by id
        activate DB
        DB --> BE: Message exists
        deactivate DB

        alt Tin nhắn tồn tại
          BE -> DB: SELECT reaction\nWHERE user + message + emoji
          activate DB
          DB --> BE: exists / not exists
          deactivate DB

          alt Đã có reaction (toggle off)
            BE -> DB: DELETE reaction
            activate DB
            DB --> BE: Deleted
            deactivate DB
          else Chưa có (thêm mới)
            BE -> DB: INSERT reaction
            activate DB
            DB --> BE: Inserted
            deactivate DB
          end

          BE -> WS: SEND reactions update\n{ messageId, reactions[] }
          activate WS
          WS --> FE: Cập nhật reactions
          deactivate WS
          BE --> FE: { reactions }
          deactivate BE
          FE --> U: Cập nhật reaction trên UI
        else Tin nhắn không tồn tại
          DB --> BE: null
          deactivate DB
          BE --> FE: 404 Not Found
          deactivate BE
          FE --> U: "Tin nhắn không tồn tại"
        end
        deactivate FE
        @enduml
    \end{plantuml}
    \caption{Sequence diagram --- Thả cảm xúc cho tin nhắn}
\end{figure}

\subsubsection{UC12 --- Trò chuyện với AI}

\renewcommand{\arraystretch}{1.4}
\begin{xltabular}{\textwidth}{|X|X|}
\hline
\multicolumn{2}{|p{0.97\textwidth}|}{\textbf{Tên usecase:} Trò chuyện với AI.} \\
\hline
\multicolumn{2}{|p{0.97\textwidth}|}{\textbf{Actor:} Người dùng.} \\
\hline
\multicolumn{2}{|p{0.97\textwidth}|}{\textbf{Mô tả:} Cho phép người dùng trò chuyện với trợ lý AI tích hợp trong ứng dụng.} \\
\hline
\multicolumn{2}{|p{0.97\textwidth}|}{\textbf{Tiền điều kiện:} Người dùng đã đăng nhập, hệ thống AI (Gemini/Groq) khả dụng.} \\
\hline
\multicolumn{2}{|p{0.97\textwidth}|}{\textbf{Hậu điều kiện:}} \\
\hline
\multicolumn{2}{|p{0.97\textwidth}|}{\quad $+$ \textbf{Thành công:} Câu trả lời của AI được hiển thị trong luồng chat AI.} \\
\hline
\multicolumn{2}{|p{0.97\textwidth}|}{\quad $+$ \textbf{Thất bại:} Hiển thị thông báo lỗi (AI không phản hồi, API key lỗi, nội dung quá dài).} \\
\hline
\multicolumn{2}{|c|}{\textbf{Luồng sự kiện chính (Basic flows):}} \\
\hline
\multicolumn{1}{|c|}{\textbf{Actor}} & \multicolumn{1}{c|}{\textbf{System}} \\
\hline
1. Người dùng mở cửa sổ chat AI từ ứng dụng. & \\
\hline
2. Nhập câu hỏi hoặc yêu cầu. & \\
\hline
& 3. Nhận yêu cầu POST /api/v1/ai/chat với nội dung tin nhắn. \\
\hline
& 4. Gửi request đến Gemini API (hoặc Groq AI với model llama-3.1-8b-instant). \\
\hline
& 5. Nhận phản hồi dạng stream từ AI. \\
\hline
& 6. Trả về phản hồi cho frontend. \\
\hline
7. Xem câu trả lời của AI trong luồng chat. & \\
\hline
\multicolumn{2}{|p{0.97\textwidth}|}{\textbf{Luồng sự kiện phụ (Alternative Flows):}} \\
\hline
\multicolumn{2}{|p{0.97\textwidth}|}{4.1 AI service không phản hồi: Backend trả về lỗi 503 Service Unavailable.} \\
\hline
\multicolumn{2}{|p{0.97\textwidth}|}{4.2 API key không hợp lệ hoặc hết hạn: Backend trả về lỗi 500 Internal Server Error.} \\
\hline
\multicolumn{2}{|p{0.97\textwidth}|}{2.1 Nội dung yêu cầu quá dài: Frontend giới hạn độ dài tin nhắn, hiển thị cảnh báo.} \\
\hline
\end{xltabular}
\renewcommand{\arraystretch}{1}

\begin{figure}[H]
    \centering
    \begin{plantuml}
        @startuml
        !theme plain
        |User|
        start
        :Mở AI Chat;
        :Nhập câu hỏi;
        :Nhấn Gửi;
        |Frontend|
        if (Nội dung trống?) then (Có - ngoại lệ)
          :Hiển thị "Vui lòng nhập câu hỏi";
          |User|
          stop
        else (Không)
          :POST /api/v1/ai/chat;
          |Backend|
          :Gửi request đến AI API;
          |AI API (Gemini/Groq)|
          :Xử lý prompt;
          if (Phản hồi thành công?) then (Có)
            |Backend|
            :Nhận phản hồi;
            |Frontend|
            :Hiển thị câu trả lời;
            |User|
            :Đọc phản hồi;
            stop
          else (Không - ngoại lệ)
            |Backend|
            |Frontend|
            :Hiển thị "AI không phản hồi\nVui lòng thử lại";
            |User|
            stop
          endif
        endif
        stop
        @enduml
    \end{plantuml}
    \caption{Activity diagram --- Trò chuyện với AI}
\end{figure}

\begin{figure}[H]
    \centering
    \begin{plantuml}
        @startuml
        !theme plain
        actor "User" as U
        participant "Frontend" as FE
        participant "Backend" as BE
        participant "AI API\n(Gemini/Groq)" as AI

        U -> FE: Nhập câu hỏi
        activate FE
        FE -> BE: POST /api/v1/ai/chat\n{ message: "..." }
        activate BE

        BE -> AI: Gửi request\n(prompt + history)
        activate AI

        alt AI phản hồi thành công
          AI --> BE: { reply: "..." }
          deactivate AI
          BE --> FE: { reply: "..." }
          deactivate BE
          FE --> U: Hiển thị câu trả lời
          deactivate FE
        else AI lỗi hoặc timeout
          AI --> BE: Error/Timeout
          deactivate AI
          BE --> FE: 503 Service Unavailable
          deactivate BE
          FE --> U: "AI hiện không khả dụng"
          deactivate FE
        end
        @enduml
    \end{plantuml}
    \caption{Sequence diagram --- Trò chuyện với AI}
\end{figure}

\subsubsection{UC13 --- Chặn người dùng}

\renewcommand{\arraystretch}{1.4}
\begin{xltabular}{\textwidth}{|X|X|}
\hline
\multicolumn{2}{|p{0.97\textwidth}|}{\textbf{Tên usecase:} Chặn người dùng.} \\
\hline
\multicolumn{2}{|p{0.97\textwidth}|}{\textbf{Actor:} Người dùng.} \\
\hline
\multicolumn{2}{|p{0.97\textwidth}|}{\textbf{Mô tả:} Cho phép người dùng chặn một người dùng khác để không nhận tin nhắn hoặc lời mời kết bạn từ họ.} \\
\hline
\multicolumn{2}{|p{0.97\textwidth}|}{\textbf{Tiền điều kiện:} Người dùng đã đăng nhập, không phải chặn chính mình.} \\
\hline
\multicolumn{2}{|p{0.97\textwidth}|}{\textbf{Hậu điều kiện:}} \\
\hline
\multicolumn{2}{|p{0.97\textwidth}|}{\quad $+$ \textbf{Thành công:} Người dùng B bị chặn, hai bên không thể gửi tin nhắn cho nhau.} \\
\hline
\multicolumn{2}{|p{0.97\textwidth}|}{\quad $+$ \textbf{Thất bại:} Đã chặn trước đó (409 Conflict), không thể chặn chính mình.} \\
\hline
\multicolumn{2}{|c|}{\textbf{Luồng sự kiện chính (Basic flows):}} \\
\hline
\multicolumn{1}{|c|}{\textbf{Actor}} & \multicolumn{1}{c|}{\textbf{System}} \\
\hline
1. Người dùng mở profile của người muốn chặn. & \\
\hline
2. Người dùng chọn "Chặn người dùng" từ menu hành động. & \\
\hline
& 3. Frontend hiển thị hộp thoại xác nhận chặn. \\
\hline
4. Người dùng xác nhận chặn. & \\
\hline
& 5. Frontend gửi POST /api/v1/block/\{userId\} đến Backend. \\
\hline
& 6. Backend kiểm tra chưa chặn trước đó, thêm record vào bảng block. \\
\hline
& 7. Backend trả về thông báo thành công, Frontend cập nhật UI. \\
\hline
\multicolumn{2}{|p{0.97\textwidth}|}{\textbf{Luồng sự kiện phụ (Alternative Flows):}} \\
\hline
\multicolumn{2}{|p{0.97\textwidth}|}{6.1 Đã chặn trước đó: Backend trả về lỗi 409 Conflict, Frontend hiển thị "Bạn đã chặn người dùng này".} \\
\hline
\multicolumn{2}{|p{0.97\textwidth}|}{4.1 Không thể chặn chính mình: Frontend ẩn nút chặn trên profile của bản thân.} \\
\hline
\caption{Mô tả chức năng UC13 --- Chặn người dùng}
\end{xltabular}
\renewcommand{\arraystretch}{1}

\begin{figure}[H]
    \centering
    \begin{plantuml}
        @startuml
        !theme plain
        |User|
        start
        :Mở profile người dùng B;
        :Chọn "Chặn người dùng";
        |Frontend|
        :Hiển thị xác nhận;
        |User|
        :Xác nhận chặn;
        |Frontend|
        :POST /block/{userId};
        |Backend|
        if (Đã chặn trước đó?) then (Có - ngoại lệ)
          |Frontend|
          :Hiển thị "Đã chặn người dùng này";
          |User|
          stop
        else (Chưa)
          :INSERT block record;
          |Database|
          :Lưu block;
          |Backend|
          |Frontend|
          :Hiển thị "Đã chặn thành công";
          |User|
          :B không thể nhắn tin cho A;
          stop
        endif
        stop
        @enduml
    \end{plantuml}
    \caption{Activity diagram --- Chặn người dùng}
\end{figure}

\begin{figure}[H]
    \centering
    \begin{plantuml}
        @startuml
        !theme plain
        actor "User A" as U
        participant "Frontend" as FE
        participant "Backend" as BE
        database "MySQL" as DB

        U -> FE: Mở profile user B
        activate FE
        U -> FE: Chọn "Chặn" + xác nhận
        FE -> BE: POST /api/v1/block/{userId}
        activate BE

        BE -> DB: SELECT block\nWHERE userId=A AND blockedId=B
        activate DB
        DB --> BE: null (chưa chặn)
        deactivate DB

        alt Chưa chặn
          BE -> DB: INSERT block
          activate DB
          DB --> BE: OK
          deactivate DB
          BE --> FE: { success }
          deactivate BE
          FE --> U: "Đã chặn người dùng"
        else Đã chặn trước đó
          DB --> BE: { block record }
          deactivate DB
          BE --> FE: 409 Conflict
          deactivate BE
          FE --> U: "Bạn đã chặn người dùng này"
        end
        deactivate FE
        @enduml
    \end{plantuml}
    \caption{Sequence diagram --- Chặn người dùng}
\end{figure}

\subsubsection{UC14 --- Báo cáo người dùng}

\renewcommand{\arraystretch}{1.4}
\begin{xltabular}{\textwidth}{|X|X|}
\hline
\multicolumn{2}{|p{0.97\textwidth}|}{\textbf{Tên usecase:} Báo cáo người dùng.} \\
\hline
\multicolumn{2}{|p{0.97\textwidth}|}{\textbf{Actor:} Người dùng.} \\
\hline
\multicolumn{2}{|p{0.97\textwidth}|}{\textbf{Mô tả:} Cho phép người dùng báo cáo người dùng khác vi phạm điều khoản sử dụng.} \\
\hline
\multicolumn{2}{|p{0.97\textwidth}|}{\textbf{Tiền điều kiện:} Người dùng đã đăng nhập, không phải báo cáo chính mình.} \\
\hline
\multicolumn{2}{|p{0.97\textwidth}|}{\textbf{Hậu điều kiện:}} \\
\hline
\multicolumn{2}{|p{0.97\textwidth}|}{\quad $+$ \textbf{Thành công:} Báo cáo được tạo với trạng thái PENDING, admin có thể xem và xử lý.} \\
\hline
\multicolumn{2}{|p{0.97\textwidth}|}{\quad $+$ \textbf{Thất bại:} Báo cáo trùng lặp (đã có report PENDING), thiếu lý do báo cáo.} \\
\hline
\multicolumn{2}{|c|}{\textbf{Luồng sự kiện chính (Basic flows):}} \\
\hline
\multicolumn{1}{|c|}{\textbf{Actor}} & \multicolumn{1}{c|}{\textbf{System}} \\
\hline
1. Người dùng mở profile người muốn báo cáo. & \\
\hline
2. Người dùng chọn "Báo cáo" từ menu hành động. & \\
\hline
& 3. Frontend hiển thị form báo cáo (lý do, mô tả, bằng chứng). \\
\hline
4. Người dùng điền lý do, đính kèm bằng chứng (nếu có) và gửi. & \\
\hline
& 5. Frontend gửi POST /api/v1/report/\{userId\} với lý do và file đính kèm. \\
\hline
& 6. Backend kiểm tra trùng lặp, tạo report với trạng thái PENDING. \\
\hline
& 7. Backend trả về thông báo thành công. \\
\hline
\multicolumn{2}{|p{0.97\textwidth}|}{\textbf{Luồng sự kiện phụ (Alternative Flows):}} \\
\hline
\multicolumn{2}{|p{0.97\textwidth}|}{6.1 Báo cáo trùng lặp: Backend từ chối nếu đã có report PENDING cho cùng cặp user.} \\
\hline
\multicolumn{2}{|p{0.97\textwidth}|}{4.1 Không thể báo cáo chính mình: Frontend ẩn nút báo cáo trên profile bản thân.} \\
\hline
\multicolumn{2}{|p{0.97\textwidth}|}{4.2 Thiếu lý do báo cáo: Backend yêu cầu lý do bắt buộc, trả về lỗi 400 Bad Request.} \\
\hline
\caption{Mô tả chức năng UC14 --- Báo cáo người dùng}
\end{xltabular}
\renewcommand{\arraystretch}{1}

\begin{figure}[H]
    \centering
    \begin{plantuml}
        @startuml
        !theme plain
        |User|
        start
        :Mở profile người vi phạm;
        :Chọn "Báo cáo";
        |Frontend|
        :Hiển thị form báo cáo;
        |User|
        :Nhập lý do;
        :Đính kèm bằng chứng (nếu có);
        :Gửi báo cáo;
        |Frontend|
        :POST /report/{userId};
        |Backend|
        if (Đã có report PENDING?) then (Có - ngoại lệ)
          |Frontend|
          :Hiển thị "Bạn đã báo cáo\nngười dùng này";
          |User|
          stop
        else (Chưa)
          if (Lý do trống?) then (Có - ngoại lệ)
            |Frontend|
            :Hiển thị "Vui lòng nhập lý do";
            |User|
            stop
          else (Không)
            :INSERT report (PENDING);
            |Database|
            :Lưu báo cáo;
            |Backend|
            |Frontend|
            :Hiển thị "Đã gửi báo cáo\nChờ admin xử lý";
            |User|
            stop
          endif
        endif
        stop
        @enduml
    \end{plantuml}
    \caption{Activity diagram --- Báo cáo người dùng}
\end{figure}

\begin{figure}[H]
    \centering
    \begin{plantuml}
        @startuml
        !theme plain
        actor "Người báo cáo" as R
        actor "Admin" as A
        participant "Frontend" as FE
        participant "Backend" as BE
        database "MySQL" as DB

        R -> FE: Chọn "Báo cáo" + nhập lý do
        activate FE
        FE -> BE: POST /api/v1/report/{userId}\n{ reason, evidence }
        activate BE

        BE -> DB: INSERT report (PENDING)
        activate DB
        DB --> BE: report saved
        deactivate DB
        BE --> FE: { success, reportId }
        deactivate BE
        FE --> R: "Đã gửi báo cáo thành công"
        deactivate FE

        == Admin xử lý sau ==

        A -> BE: GET /admin/reports\n?status=PENDING
        activate BE
        BE -> DB: SELECT reports (PENDING)
        activate DB
        DB --> BE: Danh sách báo cáo
        deactivate DB
        BE --> A: [{ reporter, target, reason }]
        deactivate BE

        A -> BE: PATCH /admin/reports/{id}/resolve\n{ action: BAN | WARN }
        activate BE
        alt Ban user
          BE -> DB: UPDATE report = RESOLVED
          BE -> DB: UPDATE user = BANNED
          activate DB
          DB --> BE: Done
          deactivate DB
        else Chỉ cảnh cáo
          BE -> DB: UPDATE report = RESOLVED
          activate DB
          DB --> BE: Done
          deactivate DB
        end
        BE --> A: { resolved }
        deactivate BE
        @enduml
    \end{plantuml}
    \caption{Sequence diagram --- Báo cáo và xử lý vi phạm}
\end{figure}

\subsubsection{UC15 --- Xử lý báo cáo (Admin)}

\renewcommand{\arraystretch}{1.4}
\begin{xltabular}{\textwidth}{|X|X|}
\hline
\multicolumn{2}{|p{0.97\textwidth}|}{\textbf{Tên usecase:} Xử lý báo cáo (Admin).} \\
\hline
\multicolumn{2}{|p{0.97\textwidth}|}{\textbf{Actor:} Quản trị viên (Admin).} \\
\hline
\multicolumn{2}{|p{0.97\textwidth}|}{\textbf{Mô tả:} Quản trị viên xem danh sách báo cáo và xử lý vi phạm bằng cách cảnh cáo hoặc khóa người dùng.} \\
\hline
\multicolumn{2}{|p{0.97\textwidth}|}{\textbf{Tiền điều kiện:} Admin đã đăng nhập với quyền ADMIN.} \\
\hline
\multicolumn{2}{|p{0.97\textwidth}|}{\textbf{Hậu điều kiện:}} \\
\hline
\multicolumn{2}{|p{0.97\textwidth}|}{\quad $+$ \textbf{Thành công:} Báo cáo được xử lý (RESOLVED), người dùng vi phạm bị cảnh cáo hoặc khóa.} \\
\hline
\multicolumn{2}{|p{0.97\textwidth}|}{\quad $+$ \textbf{Thất bại:} Báo cáo không tồn tại (404), Admin không có quyền (403).} \\
\hline
\multicolumn{2}{|c|}{\textbf{Luồng sự kiện chính (Basic flows):}} \\
\hline
\multicolumn{1}{|c|}{\textbf{Actor}} & \multicolumn{1}{c|}{\textbf{System}} \\
\hline
1. Admin mở trang quản trị báo cáo. & \\
\hline
& 2. Frontend gửi GET /api/v1/admin/reports?status=PENDING. \\
\hline
& 3. Backend truy vấn danh sách báo cáo PENDING từ database. \\
\hline
& 4. Frontend hiển thị danh sách báo cáo (người báo cáo, người bị báo cáo, lý do, bằng chứng). \\
\hline
5. Admin xem chi tiết từng báo cáo và chọn hành động xử lý. & \\
\hline
& 6. Frontend gửi PATCH /api/v1/admin/reports/\{id\}/resolve với kết quả xử lý (BAN/WARN). \\
\hline
& 7. Backend cập nhật report thành RESOLVED và thực hiện hành động (ban user nếu cần). \\
\hline
& 8. Frontend cập nhật danh sách báo cáo đã xử lý. \\
\hline
\multicolumn{2}{|p{0.97\textwidth}|}{\textbf{Luồng sự kiện phụ (Alternative Flows):}} \\
\hline
\multicolumn{2}{|p{0.97\textwidth}|}{2.1 Không có báo cáo nào: Backend trả về danh sách rỗng, Frontend hiển thị "Không có báo cáo chờ xử lý".} \\
\hline
\multicolumn{2}{|p{0.97\textwidth}|}{5.1 Báo cáo không tồn tại: Backend trả về lỗi 404 Not Found.} \\
\hline
\multicolumn{2}{|p{0.97\textwidth}|}{1.1 Admin không có quyền: Backend trả về lỗi 403 Forbidden.} \\
\hline
\caption{Mô tả chức năng UC15 --- Xử lý báo cáo (Admin)}
\end{xltabular}
\renewcommand{\arraystretch}{1}

\begin{figure}[H]
    \centering
    \begin{plantuml}
        @startuml
        !theme plain
        |Admin|
        start
        :Mở trang quản trị;
        |Frontend|
        :GET /admin/reports?status=PENDING;
        |Backend|
        :Lấy danh sách báo cáo;
        |Database|
        :SELECT PENDING reports;
        |Backend|
        |Frontend|
        if (Có báo cáo?) then (Có)
          :Hiển thị danh sách;
          |Admin|
          :Xem chi tiết từng báo cáo;
          :Quyết định hướng xử lý;
          if (Khóa user?) then (Có)
            |Frontend|
            :PATCH /admin/reports/{id}/resolve\n{ action: BAN };
            |Backend|
            :UPDATE report = RESOLVED;
            :UPDATE user = BANNED;
            |Database|
            :Cập nhật DB;
            |Admin|
            :User đã bị khóa;
          else (Chỉ cảnh cáo)
            |Frontend|
            :PATCH /admin/reports/{id}/resolve\n{ action: WARN };
            |Backend|
            :UPDATE report = RESOLVED;
            |Database|
            :Cập nhật DB;
            |Admin|
            :Báo cáo đã xử lý;
          endif
          stop
        else (Không - ngoại lệ)
          :Hiển thị "Không có\nbáo cáo chờ xử lý";
          |Admin|
          stop
        endif
        stop
        @enduml
    \end{plantuml}
    \caption{Activity diagram --- Xử lý báo cáo (Admin)}
\end{figure}

\begin{figure}[H]
    \centering
    \begin{plantuml}
        @startuml
        !theme plain
        actor "Admin" as A
        participant "Frontend" as FE
        participant "Backend" as BE
        database "MySQL" as DB

        A -> FE: Mở trang quản trị báo cáo
        activate FE
        FE -> BE: GET /admin/reports\n?status=PENDING
        activate BE
        BE -> DB: SELECT * FROM reports\nWHERE status = PENDING
        activate DB
        DB --> BE: [{ id, reporter, target, reason }]
        deactivate DB

        alt Có báo cáo
          BE --> FE: Danh sách báo cáo
          deactivate BE
          FE --> A: Hiển thị danh sách
          deactivate FE

          A -> FE: Chọn hành động: khóa user
          activate FE
          FE -> BE: PATCH /admin/reports/{id}/resolve\n{ action: BAN }
          activate BE
          BE -> DB: UPDATE report SET status = RESOLVED
          activate DB
          DB --> BE: OK
          deactivate DB
          BE -> DB: UPDATE user SET status = BANNED
          activate DB
          DB --> BE: OK
          deactivate DB
          BE --> FE: { resolved: true }
          deactivate BE
          FE --> A: "Đã xử lý - User bị khóa"
          deactivate FE
        else Không có báo cáo
          BE --> FE: []
          deactivate BE
          FE --> A: "Không có báo cáo chờ xử lý"
          deactivate FE
        end
        @enduml
    \end{plantuml}
    \caption{Sequence diagram --- Xử lý báo cáo (Admin)}
\end{figure}

\subsubsection{UC16 --- Xem thống kê hệ thống (Admin)}

\renewcommand{\arraystretch}{1.4}
\begin{xltabular}{\textwidth}{|X|X|}
\hline
\multicolumn{2}{|p{0.97\textwidth}|}{\textbf{Tên usecase:} Xem thống kê hệ thống (Admin).} \\
\hline
\multicolumn{2}{|p{0.97\textwidth}|}{\textbf{Actor:} Quản trị viên (Admin).} \\
\hline
\multicolumn{2}{|p{0.97\textwidth}|}{\textbf{Mô tả:} Quản trị viên xem thống kê tổng quan về hệ thống bao gồm số lượng người dùng, tin nhắn, nhóm, bạn bè.} \\
\hline
\multicolumn{2}{|p{0.97\textwidth}|}{\textbf{Tiền điều kiện:} Admin đã đăng nhập với quyền ADMIN.} \\
\hline
\multicolumn{2}{|p{0.97\textwidth}|}{\textbf{Hậu điều kiện:}} \\
\hline
\multicolumn{2}{|p{0.97\textwidth}|}{\quad $+$ \textbf{Thành công:} Dashboard hiển thị các chỉ số thống kê và biểu đồ.} \\
\hline
\multicolumn{2}{|p{0.97\textwidth}|}{\quad $+$ \textbf{Thất bại:} Không có dữ liệu (trả về giá trị 0), Admin không có quyền (403).} \\
\hline
\multicolumn{2}{|c|}{\textbf{Luồng sự kiện chính (Basic flows):}} \\
\hline
\multicolumn{1}{|c|}{\textbf{Actor}} & \multicolumn{1}{c|}{\textbf{System}} \\
\hline
1. Admin mở dashboard quản trị. & \\
\hline
& 2. Frontend gửi GET /api/v1/admin/stats. \\
\hline
& 3. Backend tổng hợp dữ liệu từ nhiều bảng (user, message, group, friend\_request). \\
\hline
& 4. Backend trả về các chỉ số: tổng số user, tổng số tin nhắn, số nhóm, số bạn bè. \\
\hline
& 5. Frontend hiển thị dashboard với biểu đồ và số liệu. \\
\hline
\multicolumn{2}{|p{0.97\textwidth}|}{\textbf{Luồng sự kiện phụ (Alternative Flows):}} \\
\hline
\multicolumn{2}{|p{0.97\textwidth}|}{3.1 Không có dữ liệu: Backend trả về các giá trị 0 cho tất cả chỉ số.} \\
\hline
\multicolumn{2}{|p{0.97\textwidth}|}{1.1 Admin không có quyền: Backend trả về lỗi 403 Forbidden.} \\
\hline
\caption{Mô tả chức năng UC16 --- Xem thống kê hệ thống (Admin)}
\end{xltabular}
\renewcommand{\arraystretch}{1}

\begin{figure}[H]
    \centering
    \begin{plantuml}
        @startuml
        !theme plain
        |Admin|
        start
        :Mở dashboard quản trị;
        |Frontend|
        :GET /api/v1/admin/stats;
        |Backend|
        :Tổng hợp dữ liệu;
        fork
          |Database|
          :SELECT COUNT(*) FROM users;
          |Backend|
        fork again
          |Database|
          :SELECT COUNT(*) FROM messages;
          |Backend|
        fork again
          |Database|
          :SELECT COUNT(*) FROM groups;
          |Backend|
        fork again
          |Database|
          :SELECT COUNT(*) FROM friend_requests;
          |Backend|
        end fork
        |Backend|
        :Tổng hợp kết quả;
        |Frontend|
        :Hiển thị biểu đồ + số liệu;
        |Admin|
        :Xem thống kê;
        stop
        @enduml
    \end{plantuml}
    \caption{Activity diagram --- Xem thống kê hệ thống (Admin)}
\end{figure}

\begin{figure}[H]
    \centering
    \begin{plantuml}
        @startuml
        !theme plain
        actor "Admin" as A
        participant "Frontend" as FE
        participant "Backend" as BE
        database "MySQL" as DB

        A -> FE: Mở dashboard
        activate FE
        FE -> BE: GET /api/v1/admin/stats
        activate BE

        note over BE: Truy vấn song song các bảng
        BE -> DB: SELECT COUNT(*) FROM users
        activate DB
        DB --> BE: totalUsers
        deactivate DB

        BE -> DB: SELECT COUNT(*) FROM messages
        activate DB
        DB --> BE: totalMessages
        deactivate DB

        BE -> DB: SELECT COUNT(*) FROM groups
        activate DB
        DB --> BE: totalGroups
        deactivate DB

        BE -> DB: SELECT COUNT(*) FROM friend_requests
        activate DB
        DB --> BE: totalFriendRequests
        deactivate DB

        BE --> FE: { totalUsers, totalMessages,\ntotalGroups, totalFriendRequests }
        deactivate BE
        FE --> A: Hiển thị biểu đồ và số liệu
        deactivate FE
        @enduml
    \end{plantuml}
    \caption{Sequence diagram --- Xem thống kê hệ thống (Admin)}
\end{figure}

\subsubsection{UC17 --- Ghi log cuộc gọi}

\renewcommand{\arraystretch}{1.4}
\begin{xltabular}{\textwidth}{|X|X|}
\hline
\multicolumn{2}{|p{0.97\textwidth}|}{\textbf{Tên usecase:} Ghi log cuộc gọi.} \\
\hline
\multicolumn{2}{|p{0.97\textwidth}|}{\textbf{Actor:} Người dùng, Hệ thống.} \\
\hline
\multicolumn{2}{|p{0.97\textwidth}|}{\textbf{Mô tả:} Hệ thống tự động ghi lại thông tin cuộc gọi thoại/video giữa các người dùng sau khi cuộc gọi kết thúc.} \\
\hline
\multicolumn{2}{|p{0.97\textwidth}|}{\textbf{Tiền điều kiện:} Người dùng đã đăng nhập, WebRTC đã được thiết lập.} \\
\hline
\multicolumn{2}{|p{0.97\textwidth}|}{\textbf{Hậu điều kiện:}} \\
\hline
\multicolumn{2}{|p{0.97\textwidth}|}{\quad $+$ \textbf{Thành công:} Log cuộc gọi được lưu vào database, hiển thị trong lịch sử chat.} \\
\hline
\multicolumn{2}{|p{0.97\textwidth}|}{\quad $+$ \textbf{Thất bại:} Cuộc gọi không kết nối được (ghi log MISSED), mất kết nối giữa chừng (ghi thời lượng thực tế).} \\
\hline
\multicolumn{2}{|c|}{\textbf{Luồng sự kiện chính (Basic flows):}} \\
\hline
\multicolumn{1}{|c|}{\textbf{Actor}} & \multicolumn{1}{c|}{\textbf{System}} \\
\hline
1. Người dùng A thực hiện cuộc gọi đến người dùng B. & \\
\hline
& 2. Hệ thống WebRTC thiết lập kết nối giữa hai bên. \\
\hline
3. Người dùng B chấp nhận cuộc gọi. & \\
\hline
& 4. Frontend thiết lập WebRTC PeerConnection, trao đổi media stream. \\
\hline
5. Một trong hai bên kết thúc cuộc gọi. & \\
\hline
& 6. Frontend gửi POST /api/v1/calls với thông tin cuộc gọi. \\
\hline
& 7. Backend ghi log cuộc gọi (thời gian bắt đầu, kết thúc, thời lượng, loại, trạng thái). \\
\hline
& 8. Frontend hiển thị lịch sử cuộc gọi trong danh sách chat. \\
\hline
\multicolumn{2}{|p{0.97\textwidth}|}{\textbf{Luồng sự kiện phụ (Alternative Flows):}} \\
\hline
\multicolumn{2}{|p{0.97\textwidth}|}{3.1 Cuộc gọi không kết nối được (B từ chối hoặc không trả lời): Ghi log với trạng thái MISSED.} \\
\hline
\multicolumn{2}{|p{0.97\textwidth}|}{7.1 Lỗi database: Log không được lưu nhưng không ảnh hưởng đến cuộc gọi.} \\
\hline
\caption{Mô tả chức năng UC17 --- Ghi log cuộc gọi}
\end{xltabular}
\renewcommand{\arraystretch}{1}

\begin{figure}[H]
    \centering
    \begin{plantuml}
        @startuml
        !theme plain
        |User A|
        start
        :Chọn gọi cho User B;
        :Chọn loại (thoại/video);
        |Frontend|
        :Tạo WebRTC offer;
        |User B|
        if (B chấp nhận?) then (Có)
          :Chấp nhận cuộc gọi;
          |Frontend|
          :Thiết lập WebRTC PeerConnection;
          :Trao đổi media stream;
          note right: Cuộc gọi đang diễn ra
          |User A|
          :Kết thúc cuộc gọi;
          |Frontend|
          :POST /api/v1/calls;
          |Backend|
          :Ghi log cuộc gọi;
          |Database|
          :INSERT call log\n(thời lượng, loại, OK);
          |User A|
          :Xem lịch sử cuộc gọi;
          |User B|
          stop
        else (Không - ngoại lệ)
          |Backend|
          :Ghi log MISSED;
          |Database|
          :INSERT call log\n(status: MISSED);
          |User A|
          :Thấy "Cuộc gọi nhỡ";
          stop
        endif
        stop
        @enduml
    \end{plantuml}
    \caption{Activity diagram --- Ghi log cuộc gọi}
\end{figure}

\begin{figure}[H]
    \centering
    \begin{plantuml}
        @startuml
        !theme plain
        actor "User A" as A
        participant "Frontend A" as FEA
        participant "WebRTC" as RTC
        actor "User B" as B
        participant "Frontend B" as FEB
        participant "Backend" as BE
        database "MySQL" as DB

        A -> FEA: Gọi cho B (thoại/video)
        activate FEA
        FEA -> RTC: Tạo offer
        RTC -> FEB: Thông báo cuộc gọi đến
        FEB --> B: Hiển thị màn hình cuộc gọi

        alt B chấp nhận
          B -> FEB: Chấp nhận
          FEB -> RTC: Answer
          RTC --> FEA: Kết nối established
          note over RTC, FEB: Trao đổi media stream (WebRTC)

          A -> FEA: Kết thúc cuộc gọi
          deactivate FEA
          FEA -> BE: POST /api/v1/calls\n{ type, duration, status: OK }
          activate BE
          BE -> DB: INSERT call log
          activate DB
          DB --> BE: Saved
          deactivate DB
          BE --> FEA: Log saved
          deactivate BE
        else B từ chối/không trả lời
          B -> FEB: Từ chối / timeout
          FEA -> BE: POST /api/v1/calls\n{ status: MISSED }
          activate BE
          BE -> DB: INSERT call log (MISSED)
          activate DB
          DB --> BE: Saved
          deactivate DB
          BE --> FEA: Log saved
          deactivate BE
        end
        @enduml
    \end{plantuml}
    \caption{Sequence diagram --- Ghi log cuộc gọi}
\end{figure}

\section{Class diagram}

\begin{figure}[H]
    \centering
    \begin{plantuml}
        @startuml
        !theme plain
        skinparam classAttributeIconSize 0
        skinparam defaultFontSize 11

        class BaseAuditingEntity {
            +createdDate: LocalDateTime
            +lastModifiedDate: LocalDateTime
        }

        class User {
            +id: UUID
            +firstName: String
            +lastName: String
            +email: String
            +password: String
            +role: String
            +avatarUrl: String
            +banned: boolean
            +banReason: String
            +banUntil: LocalDateTime
            +bannedAt: LocalDateTime
            +emailVerified: boolean
            +online: boolean
            +lastSeen: LocalDateTime
            +tokenVersion: int
        }

        class Chat {
            +id: UUID
            +user1: User
            +user2: User
            +deletedByUser1: boolean
            +deletedByUser2: boolean
        }

        class Message {
            +id: UUID
            +content: String
            +type: MessageType
            +state: MessageState
            +fileName: String
            +deleted: boolean
            +chat: Chat
            +sender: User
        }

        class Group {
            +id: UUID
            +name: String
            +avatarUrl: String
            +description: String
            +createdBy: User
        }

        class GroupMember {
            +id: UUID
            +group: Group
            +user: User
            +admin: boolean
        }

        class GroupMessage {
            +id: UUID
            +content: String
            +type: MessageType
            +fileName: String
            +deleted: boolean
            +group: Group
            +sender: User
        }

        class FriendRequest {
            +id: UUID
            +sender: User
            +receiver: User
            +status: String
        }

        class MessageReaction {
            +id: UUID
            +message: Message
            +user: User
            +emoji: String
        }

        class GroupMessageReaction {
            +id: UUID
            +groupMessage: GroupMessage
            +user: User
            +emoji: String
        }

        class BlockedUser {
            +id: UUID
            +blocker: User
            +blocked: User
        }

        class CallSession {
            +id: UUID
            +chat: Chat
            +initiator: User
            +receiver: User
            +callType: String
            +status: String
            +durationSec: Integer
        }

        class AiMessage {
            +id: UUID
            +user: User
            +role: String
            +content: String
        }

        class GroupJoinRequest {
            +id: UUID
            +group: Group
            +requestedBy: User
            +targetUser: User
            +status: String
        }

        class Report {
            +id: Long
            +reporter: User
            +reported: User
            +reason: String
            +description: String
            +status: String
            +evidenceKeys: String
        }

        class AuditLog {
            +id: UUID
            +adminId: UUID
            +adminEmail: String
            +action: String
            +targetType: String
            +targetName: String
            +details: String
        }

        BaseAuditingEntity <|-- User
        BaseAuditingEntity <|-- Chat
        BaseAuditingEntity <|-- Message
        BaseAuditingEntity <|-- Group
        BaseAuditingEntity <|-- GroupMember
        BaseAuditingEntity <|-- GroupMessage
        BaseAuditingEntity <|-- FriendRequest
        BaseAuditingEntity <|-- MessageReaction
        BaseAuditingEntity <|-- GroupMessageReaction
        BaseAuditingEntity <|-- CallSession
        BaseAuditingEntity <|-- AiMessage
        BaseAuditingEntity <|-- GroupJoinRequest

        Chat "1" *-- "2" User : contains
        Message "*" *-- "1" Chat : belongs to
        Message "*" *-- "1" User : sent by
        GroupMember "*" -- "1" Group
        GroupMember "*" -- "1" User
        GroupMessage "*" -- "1" Group
        GroupMessage "*" -- "1" User
        FriendRequest "*" -- "1" User : sender
        FriendRequest "*" -- "1" User : receiver
        MessageReaction "*" -- "1" Message
        MessageReaction "*" -- "1" User
        BlockedUser "*" -- "1" User : blocker
        BlockedUser "*" -- "1" User : blocked
        CallSession "*" -- "1" Chat
        AiMessage "*" -- "1" User
        GroupJoinRequest "*" -- "1" Group
        GroupJoinRequest "*" -- "1" User : requestedBy
        GroupJoinRequest "*" -- "1" User : targetUser
        @enduml
    \end{plantuml}
    \caption{Class diagram --- Mô hình lớp tổng quát}
\end{figure}

\section{Architecture diagram}

\begin{figure}[H]
    \centering
    \begin{plantuml}
        @startuml
        !theme plain
        skinparam componentStyle rectangle

        package "Frontend (Vercel)" {
            [React App] as FE
            [Zustand Store] as Store
            [WebSocket Client] as WSClient
            [Axios API Layer] as API
        }

        package "Backend (Railway)" {
            [Spring Boot REST API] as REST
            [Spring Security + JWT] as Security
            [WebSocket Broker] as WSBroker
            [Message Service] as MsgSvc
            [User Service] as UserSvc
            [Group Service] as GroupSvc
            [AI Service] as AISvc
            [File Storage Service] as FileSvc
            [Report Service] as ReportSvc
        }

        database "MySQL (Railway)" as DB
        database "Redis (Railway)" as Redis
        node "Cloudflare R2" as R2
        node "Groq API" as Groq
        node "Brevo API" as Brevo

        FE --> API : HTTP REST
        FE --> WSClient : STOMP over WebSocket
        API --> REST : HTTP calls
        WSClient --> WSBroker : STOMP
        REST --> Security : JWT filter
        REST --> MsgSvc
        REST --> UserSvc
        REST --> GroupSvc
        REST --> AISvc
        REST --> FileSvc
        REST --> ReportSvc
        MsgSvc --> DB : JPA
        UserSvc --> DB
        GroupSvc --> DB
        FileSvc --> R2 : S3 SDK
        FileSvc --> DB
        AISvc --> Groq : HTTP
        UserSvc --> Brevo : Email
        REST --> Redis : Cache
        WSBroker --> Redis : Pub/Sub
        @enduml
    \end{plantuml}
    \caption{Architecture diagram --- Kiến trúc tổng thể hệ thống}
\end{figure}

\section{Deployment diagram}

\begin{figure}[H]
    \centering
    \begin{plantuml}
        @startuml
        !theme plain
        skinparam componentStyle rectangle

        node "User's Browser" {
            [React SPA]
        }

        node "Vercel" {
            [Vercel CDN]
            artifact "index.html + JS bundle"
        }

        node "Railway (asia-southeast1)" {
            artifact "Docker Container" {
                [Spring Boot App :8080]
            }
            database "MySQL Addon :3306" {
                [zalo\_clone DB]
            }
            database "Redis Addon :6379" {
                [Cache + Session]
            }
        }

        node "Cloudflare R2" {
            [File Storage Bucket]
        }

        node "Groq Cloud" {
            [LLaMA 3.1 8B API]
        }

        node "Brevo" {
            [Email Service API]
        }

        "User's Browser" --> "Vercel" : HTTPS (tải frontend)
        "User's Browser" --> "Railway (asia-southeast1)" : HTTPS (API calls)
        "User's Browser" --> "Railway (asia-southeast1)" : WSS (WebSocket)
        "Railway (asia-southeast1)" --> "Cloudflare R2" : S3 API
        "Railway (asia-southeast1)" --> "Groq Cloud" : OpenAI-compatible API
        "Railway (asia-southeast1)" --> "Brevo" : REST API

        note "Backend: Spring Boot 3.4.2\nJava 21, Docker\nFrontend: React 19, Vite" as N
        N -- "Railway (asia-southeast1)"
        @enduml
    \end{plantuml}
    \caption{Deployment diagram --- Mô hình triển khai hệ thống}
\end{figure}

% #############################################################################
% CHƯƠNG 4: HIỆN THỰC
% #############################################################################
\chapter{HIỆN THỰC}

\section{Cấu hình phần cứng, phần mềm}

\subsection{Backend Server --- Railway (PaaS)}

\begin{itemize}
    \item \textbf{Nền tảng}: Railway, vùng Asia Southeast 1 (Singapore).
    \item \textbf{Container}: Docker, JDK 21 (eclipse-temurin:21-jre), JVM heap 256-512MB.
    \item \textbf{CPU/RAM}: Phân bổ động theo nhu cầu, tối đa 512MB RAM.
    \item \textbf{Storage}: MySQL addon đi kèm, Cloudflare R2 cho file upload.
\end{itemize}

\subsection{Frontend --- Vercel}

\begin{itemize}
    \item \textbf{Triển khai}: Vercel CDN toàn cầu.
    \item \textbf{Build}: Vite 8, React 19, bundle tối ưu.
    \item \textbf{Domain}: \url{https://zalo-fullstack-app.vercel.app}.
\end{itemize}

\subsection{Máy phát triển (Dev Environment)}

\begin{itemize}
    \item \textbf{Hệ điều hành}: Windows 11.
    \item \textbf{CPU/RAM}: Intel Core i7 / 16GB RAM.
    \item \textbf{Java}: JDK 21.
    \item \textbf{Node.js}: Bản mới nhất (dùng cho frontend dev).
    \item \textbf{MySQL}: 8.0 local hoặc Railway MySQL addon.
\end{itemize}

\subsection{Công nghệ sử dụng}

\begin{table}[H]
    \centering
    \caption{Công nghệ và phiên bản}
    \begin{tabularx}{\textwidth}{|l|X|l|}
        \hline
        \textbf{Công nghệ} & \textbf{Mô tả} & \textbf{Phiên bản} \\
        \hline
        Java & Ngôn ngữ lập trình backend & 21 \\
        \hline
        Spring Boot & Framework backend & 3.4.2 \\
        \hline
        MySQL & Cơ sở dữ liệu quan hệ & 8.0 \\
        \hline
        React & Thư viện frontend & 19.2.4 \\
        \hline
        Vite & Công cụ build frontend & 8.0 \\
        \hline
        Zustand & Quản lý state frontend & 5.0.12 \\
        \hline
        STOMP/WebSocket & Giao thức real-time & 7.3 \\
        \hline
        JWT & Xác thực token & 0.11.5 \\
        \hline
        Flyway & Quản lý migration DB & 10 \\
        \hline
        Docker & Container hóa & 24+ \\
        \hline
    \end{tabularx}
\end{table}

\section{Giao diện hệ thống}

\subsection{Màn hình đăng nhập}

Màn hình đăng nhập cho phép người dùng nhập email và mật khẩu để đăng nhập. Hỗ trợ đăng ký tài khoản mới và quên mật khẩu. Nếu tài khoản bị khóa, hệ thống hiển thị card thông báo riêng với lý do và thời hạn.

\begin{figure}[H]
    \centering
    \begin{minipage}{0.45\textwidth}
        \centering
        \includegraphics[width=\textwidth]{login-screen.png}
        \subcaption{Web}
    \end{minipage}
    \hfill
    \begin{minipage}{0.45\textwidth}
        \centering
        \includegraphics[width=\textwidth]{login-screen-mobile.png}
        \subcaption{Mobile}
    \end{minipage}
    \caption{Màn hình đăng nhập}
\end{figure}

\begin{figure}[H]
    \centering
    \begin{minipage}{0.45\textwidth}
        \centering
        \includegraphics[width=\textwidth]{register-screen.png}
        \subcaption{Web}
    \end{minipage}
    \hfill
    \begin{minipage}{0.45\textwidth}
        \centering
        \includegraphics[width=\textwidth]{register-screen-mobile.png}
        \subcaption{Mobile}
    \end{minipage}
    \caption{Màn hình đăng ký}
\end{figure}

\subsection{Màn hình chat chính}

Giao diện chính của ứng dụng gồm:

\begin{itemize}
    \item \textbf{Sidebar bên trái}: Danh sách chat, nhóm, tìm kiếm người dùng, feed, AI chat, danh bạ. Mỗi chat hiển thị avatar, tên, tin nhắn cuối, thời gian, và số tin nhắn chưa đọc.
    \item \textbf{Ở giữa}: Cửa sổ chat chính. Tin nhắn hiển thị dạng bubble, tin nhắn của mình màu xanh (\#0068ff), tin nhắn đối phương màu trắng. Có input để gửi văn bản, đính kèm file, emoji.
    \item \textbf{Bên phải (tuỳ chọn)}: Thông tin chat/người dùng.
\end{itemize}

\begin{figure}[H]
    \centering
    \begin{minipage}{0.45\textwidth}
        \centering
        \includegraphics[width=\textwidth]{main-chat.png}
        \subcaption{Web}
    \end{minipage}
    \hfill
    \begin{minipage}{0.45\textwidth}
        \centering
        \includegraphics[width=\textwidth]{main-chat-mobile.png}
        \subcaption{Mobile}
    \end{minipage}
    \caption{Màn hình chat chính}
\end{figure}

\subsection{Màn hình nhóm chat}

Cho phép tạo nhóm với tên, mô tả, avatar. Quản lý thành viên: thêm/xóa thành viên, phân quyền admin nhóm. Tin nhắn nhóm được broadcast qua WebSocket đến tất cả thành viên.

\begin{figure}[H]
    \centering
    \begin{minipage}{0.45\textwidth}
        \centering
        \includegraphics[width=\textwidth]{group-chat.png}
        \subcaption{Web}
    \end{minipage}
    \hfill
    \begin{minipage}{0.45\textwidth}
        \centering
        \includegraphics[width=\textwidth]{group-chat-mobile.png}
        \subcaption{Mobile}
    \end{minipage}
    \caption{Màn hình nhóm chat}
\end{figure}

\subsection{Màn hình quản trị (Admin)}

Dashboard thống kê với các chỉ số: tổng người dùng, tin nhắn, nhóm; số người online, bị ban; biểu đồ tin nhắn/người dùng mới theo ngày. Quản lý danh sách người dùng, nhóm, báo cáo vi phạm.

\begin{figure}[H]
    \centering
    \includegraphics[width=\textwidth]{admin-dashboard.png}
    \caption{Màn hình quản trị (Admin)}
\end{figure}

\subsection{Tính năng AI Chat}

Giao diện chat riêng với AI bot, sử dụng Groq API (LLaMA 3.1 8B). Hệ thống hỗ trợ ba chế độ tương tác AI:

\begin{itemize}
    \item \textbf{Chat AI trực tiếp}: Người dùng nhập câu hỏi và nhận phản hồi từ AI theo thời gian thực.
    \item \textbf{AI gợi ý}: AI tự động đề xuất tin nhắn trả lời dựa trên ngữ cảnh hội thoại.
    \item \textbf{AI tóm tắt chat}: AI tóm tắt nội dung hội thoại dài thành các ý chính ngắn gọn.
\end{itemize}

\begin{figure}[H]
    \centering
    \includegraphics[width=0.8\textwidth]{ai1.jpg}
    \caption{Chat AI trực tiếp}
\end{figure}

\begin{figure}[H]
    \centering
    \includegraphics[width=0.8\textwidth]{ai2.jpg}
    \caption{AI gợi ý tin nhắn}
\end{figure}

\begin{figure}[H]
    \centering
    \includegraphics[width=0.8\textwidth]{ai3.jpg}
    \caption{AI tóm tắt chat}
\end{figure}

\section{Kế hoạch và hiện thực kiểm thử hệ thống}

\subsection{Kế hoạch kiểm thử}

\begin{itemize}
    \item \textbf{Unit Test}: Spring Boot Test với JUnit 5, Mockito. Kiểm tra service layer, mapper, exception handling.
    \item \textbf{Integration Test}: Testcontainers với MySQL container. Kiểm tra toàn bộ luồng từ controller đến database.
    \item \textbf{API Test}: Postman collection cho tất cả endpoint.
    \item \textbf{Frontend Test}: ESLint kiểm tra code quality.
\end{itemize}

\subsection{Kiểm thử hệ thống}

Các kịch bản kiểm thử chính:

\begin{longtable}{|c|p{9cm}|c|}
    \hline
    \textbf{ID} & \textbf{Kịch bản} & \textbf{Kết quả} \\
    \hline
    T01 & Đăng ký tài khoản mới với email hợp lệ & Pass \\
    \hline
    T02 & Đăng nhập với thông tin đúng & Pass \\
    \hline
    T03 & Đăng nhập với sai mật khẩu & Pass (trả về lỗi) \\
    \hline
    T04 & Gửi tin nhắn văn bản 1-1 & Pass \\
    \hline
    T05 & Gửi tin nhắn với file đính kèm & Pass \\
    \hline
    T06 & Tạo nhóm chat & Pass \\
    \hline
    T07 & Thêm/xóa thành viên nhóm & Pass \\
    \hline
    T08 & Gửi tin nhắn nhóm & Pass \\
    \hline
    T09 & Thả reaction cho tin nhắn & Pass \\
    \hline
    T10 & Kết bạn -- gửi và chấp nhận & Pass \\
    \hline
    T11 & Chặn người dùng & Pass \\
    \hline
    T12 & Báo cáo vi phạm & Pass \\
    \hline
    T13 & Admin xử lý báo cáo + ban user & Pass \\
    \hline
    T14 & AI Chat với Groq API & Pass \\
    \hline
    T15 & Refresh token tự động & Pass \\
    \hline
    T16 & Đăng nhập từ thiết bị mới -- kick thiết bị cũ & Pass \\
    \hline
    T17 & Upload file -- download file & Pass \\
    \hline
    \caption{Kết quả kiểm thử hệ thống}
\end{longtable}

% #############################################################################
% CHƯƠNG 5: KẾT LUẬN
% #############################################################################
\chapter{KẾT LUẬN}

\section{Kết quả đạt được}

Đồ án đã hoàn thành mục tiêu xây dựng ứng dụng nhắn tin Zalo Clone với các kết quả sau:

\begin{itemize}
    \item \textbf{Backend hoàn chỉnh}: 26 migration Flyway quản lý schema với 16+ bảng, 30+ REST API endpoint, hệ thống xác thực JWT với token version, WebSocket real-time cho chat 1-1 và nhóm.
    \item \textbf{Frontend hiện đại}: Giao diện chat real-time với React 19, quản lý state với Zustand, WebSocket tự động kết nối lại với singleton pattern.
    \item \textbf{Đa phương tiện}: Upload và stream hình ảnh, video, audio, file với Cloudflare R2, hỗ trợ HTTP Range request cho video streaming.
    \item \textbf{Tính năng xã hội}: Kết bạn, chặn người dùng, báo cáo vi phạm, quản lý nhóm.
    \item \textbf{AI Integration}: Tích hợp Groq API với LLaMA 3.1, lưu lịch sử chat AI.
    \item \textbf{Admin Dashboard}: Thống kê real-time, quản lý người dùng/nhóm, audit log, xử lý báo cáo.
    \item \textbf{Triển khai Cloud}: Backend trên Railway (Docker), frontend trên Vercel, file storage trên Cloudflare R2.
\end{itemize}

\section{Hạn chế của đồ án}

Bên cạnh những kết quả đạt được, đồ án còn một số hạn chế:

\begin{itemize}
    \item \textbf{Thiếu gọi điện thoại/video real-time}: Chỉ ghi nhận log call session qua WebRTC signal chưa được implement đầy đủ.
    \item \textbf{Chưa có end-to-end encryption}: Tin nhắn lưu dạng plaintext trong database.
    \item \textbf{Chưa có real-time notification ngoài WebSocket}: Thiếu push notification (FCM/APNs) khi tab trình duyệt không active.
\end{itemize}

\section{Hướng phát triển}

\begin{itemize}
    \item \textbf{Video Call / Voice Call}: Implement WebRTC real-time với TURN server.
    \item \textbf{End-to-End Encryption}: Mã hóa tin nhắn đầu cuối với Signal Protocol.
    \item \textbf{Push Notification}: Tích hợp Firebase Cloud Messaging.
    \item \textbf{Story / Status}: Tính năng đăng trạng thái tạm thời (24h).
    \item \textbf{Cải thiện hiệu năng}: Phân trang chat list, virtual scrolling cho tin nhắn.
\end{itemize}

% #############################################################################
% TÀI LIỆU THAM KHẢO
% #############################################################################
\chapter*{TÀI LIỆU THAM KHẢO}
\addcontentsline{toc}{chapter}{TÀI LIỆU THAM KHẢO}

\section*{Tài liệu Tiếng Anh}

\begin{enumerate}
    \item Spring Boot Reference Documentation. \url{https://docs.spring.io/spring-boot/docs/3.4.2/reference/} --- Truy cập 04/2026.
    \item React 19 Documentation. \url{https://react.dev/blog/2024/12/05/react-19} --- Truy cập 04/2026.
    \item STOMP Protocol Specification, Version 1.2. \url{https://stomp.github.io/stomp-specification-1.2.html} --- Truy cập 04/2026.
    \item JSON Web Token (JWT) --- RFC 7519. \url{https://datatracker.ietf.org/doc/html/rfc7519} --- Truy cập 04/2026.
    \item Flyway Documentation. \url{https://documentation.red-gate.com/fd/flyway-10} --- Truy cập 04/2026.
    \item Railway Documentation. \url{https://docs.railway.app/} --- Truy cập 04/2026.
    \item Groq API Documentation. \url{https://console.groq.com/docs} --- Truy cập 04/2026.
\end{enumerate}

\section*{Tài liệu từ Internet}

\begin{enumerate}
    \setcounter{enumi}{7}
    \item Zustand State Management. \url{https://zustand.docs.pmnd.rs/getting-started/introduction} --- Truy cập 04/2026.
    \item Tailwind CSS Documentation. \url{https://tailwindcss.com/docs} --- Truy cập 04/2026.
    \item Cloudflare R2 Documentation. \url{https://developers.cloudflare.com/r2/} --- Truy cập 04/2026.
    \item Vite Documentation. \url{https://vite.dev/guide/} --- Truy cập 04/2026.
\end{enumerate}

\end{document}
